% ============================================================================
% sec10.tex — §10 Steady-State Heat Conduction, Fourier's Law, Composite Walls,
%             Cylinders/Spheres, Contact Resistance & Critical Radius
% Target: 12 pages | Figures: fig:10:composite-wall, fig:10:cylindrical-rcrit
% ============================================================================
\binderchapter{10}{Steady-State Heat Conduction, Fourier's Law, Composite Walls, Cylinders/Spheres, Contact Resistance \& Critical Radius}{§10 Heat Conduction}{sec:10}
\scopebanner{REQUIRED}[2027 Rule 5.c.iii (All Tournament Levels: Kinds of Heat Transfer, Thermal Conductivity) \textbar\ Sources: \src{S01}, \src{S12}, \src{S13}, \src{S17}, \src{S25}]

\idxtopic{conduction!Fourier's law}
\idxtopic{thermal conductivity!definition and mechanisms}
\idxtopic{thermal resistance!electrical analogy}
\idxtopic{thermal diffusivity}

% ----------------------------------------------------------------------------
% L1: FORMULA & DECISION CARD
% ----------------------------------------------------------------------------
\begin{formulacard}[§10 Master Conduction, Thermal Resistance, Fin \& Wiedemann--Franz Decision Card \src{S12,S13,S25}]
\begin{ftab}
$\dot{Q} = k A \dfrac{T_1 - T_2}{L} = \dfrac{\Delta T}{R_{\mathrm{slab}}}$
  & $\dot{Q}$: conduction rate $(\si{\watt})$; $k$: thermal conductivity $(\si{\watt\per\meter\per\kelvin})$; $A$: cross-sectional area $\perp$ heat flow $(\si{\square\meter})$; $L$: thickness $(\si{\meter})$; $\Delta T = T_1 - T_2 > 0$ $(\si{\kelvin})$
  & \flag{1D-STEADY}
  & Plane slab/wall of constant area $A$ and homogeneous $k$ \idxformula{Fourier's law!plane slab $\dot{Q}=kA\Delta T/L$} \\
$q'' = \dfrac{\dot{Q}}{A} = -k\dfrac{\Delta T}{\Delta x}$
  & $q''$: heat flux $(\si{\watt\per\square\meter})$; minus sign indicates heat flows down the temperature gradient ($T_{\mathrm{hot}} \to T_{\mathrm{cold}}$)
  & \flag{1D-STEADY}
  & Heat flow per unit area; constant across plane series layers \idxformula{heat flux!$q''=\dot{Q}/A$} \\
$\alpha = \dfrac{k}{\rho c_p}, \quad b = \sqrt{k\rho c_p}$
  & $\alpha$: thermal diffusivity $(\si{\square\meter\per\second})$; $b$: thermal effusivity $(\si{\joule\per\square\meter\per\kelvin\per\second\tothe{1/2}})$; $\rho c_p$: volumetric heat capacity $(\si{\joule\per\cubic\meter\per\kelvin})$
  & \CONSTC
  & Speed of thermal penetration ($\delta_t \sim 2\sqrt{\alpha t}$) vs.\ contact heat sharing \idxformula{thermal diffusivity!$\alpha=k/(\rho c_p)$} \\
$R_{\mathrm{slab}} = \dfrac{L}{k A}, \quad R_{\mathrm{val}} = \dfrac{L}{k}$
  & $R_{\mathrm{slab}}$: lumped resistance $(\si{\kelvin\per\watt})$; $R_{\mathrm{val}} = A R_{\mathrm{slab}}$: unit-area $R$-value $(\si{\square\meter\kelvin\per\watt})$ [$1\,\text{RSI} = 5.678\,\text{R}_{\text{US}}$]
  & \flag{1D-STEADY}
  & Plane walls; series resistances add ($R_{\mathrm{tot}} = \sum R_i$) \idxformula{thermal resistance!plane slab $R=L/(kA)$} \\
$R_{\mathrm{cyl}} = \dfrac{\ln(r_o / r_i)}{2\pi k \ell} = \dfrac{r_o - r_i}{k A_{\mathrm{lm}}}$
  & $r_i, r_o$: inner/outer radii $(\si{\meter})$; $\ell$: axial length $(\si{\meter})$; log-mean area $A_{\mathrm{lm}} = \dfrac{2\pi\ell(r_o - r_i)}{\ln(r_o/r_i)}$ $(\si{\square\meter})$
  & \flag{1D-RADIAL}
  & Cylindrical pipe, tube, or 2027 cup sidewall ($\dot{Q} = \Delta T / R_{\mathrm{cyl}}$) \idxformula{thermal resistance!cylinder $R=\ln(r_o/r_i)/(2\pi k\ell)$} \\
$R_{\mathrm{sph}} = \dfrac{r_o - r_i}{4\pi k r_i r_o} = \dfrac{r_o - r_i}{k A_{\mathrm{gm}}}$
  & $r_i, r_o$: inner/outer radii $(\si{\meter})$; geometric-mean area $A_{\mathrm{gm}} = \sqrt{A_i A_o} = 4\pi r_i r_o$ $(\si{\square\meter})$
  & \flag{1D-RADIAL}
  & Spherical shells, Dewar bulbs, or quasi-spherical cavities \idxformula{thermal resistance!sphere $R=(r_o-r_i)/(4\pi k r_i r_o)$} \\
$R_{\mathrm{tot,ser}} = \sum_{i=1}^n R_i, \quad \dfrac{1}{R_{\mathrm{tot,par}}} = \sum_{j=1}^m \dfrac{1}{R_j}$
  & Series: common $\dot{Q}$, $\Delta T = \dot{Q} R_{\mathrm{tot,ser}}$; Parallel: common $\Delta T$, $\dot{Q}_{\mathrm{tot}} = \sum \dot{Q}_j = \Delta T / R_{\mathrm{tot,par}}$
  & \flag{NETWORK}
  & Multi-layer walls, studs/bolts (thermal bridges), combined modes \idxformula{thermal resistance!series and parallel networks} \\
$U A = \dfrac{1}{R_{\mathrm{tot}}}, \quad \dot{Q} = U A \Delta T_{\mathrm{overall}}$
  & $U$: overall heat transfer coefficient $(\si{\watt\per\square\meter\per\kelvin})$ referenced to chosen area $A$ ($U_i A_i = U_o A_o$)
  & \flag{COMPOSITE}
  & Heat exchangers, walls with indoor/outdoor convection $h_{\mathrm{in}}, h_{\mathrm{out}}$ \idxformula{overall heat transfer coefficient!$UA=1/R_{\mathrm{tot}}$} \\
$R_c = \dfrac{1}{h_c A}, \quad \Delta T_c = \dfrac{\dot{Q}}{h_c A} = \dfrac{q''}{h_c}$
  & $h_c$: interfacial contact conductance $(\si{\watt\per\square\meter\per\kelvin})$ (Lienhard Table 2.1 \src{S12}); $R_c'' = 1/h_c$ $(\si{\square\meter\kelvin\per\watt})$
  & \flag{INTERFACE}
  & Finite temperature jump across rough solid--solid joints \idxformula{contact resistance!$R_c=1/(h_c A)$} \\
$r_{\mathrm{crit,cyl}} = \dfrac{k_{\mathrm{ins}}}{h_{\mathrm{ext}}}, \quad r_{\mathrm{crit,sph}} = \dfrac{2k_{\mathrm{ins}}}{h_{\mathrm{ext}}}$
  & $r_{\mathrm{crit}}$: outer insulation radius maximizing $\dot{Q}$ (minimizing $R_{\mathrm{tot}}$), where $\mathrm{Bi}_{r_{\mathrm{crit}}} = h_{\mathrm{ext}}r_{\mathrm{crit}}/k_{\mathrm{ins}} = 1$ (cyl) or $2$ (sph)
  & \flag{RADIAL-OPT}
  & Thin wires/tubes ($r_i < r_{\mathrm{crit}}$); for 9 oz cup $r_i \gg r_{\mathrm{crit}}$ so insulation always helps \idxformula{critical radius of insulation!$r_{\mathrm{crit}}=k/h$} \\
$m = \sqrt{\dfrac{h P}{k A_c}}, \quad \eta_f = \dfrac{\tanh(m L_c)}{m L_c}$
  & $P$: fin perimeter $(\si{\meter})$; $A_c$: fin cross-section $(\si{\square\meter})$; $L_c = L + A_c/P$; $\dot{Q}_f = \eta_f h A_f \theta_b$; effectiveness $\varepsilon_f = \dot{Q}_f / (h A_c \theta_b)$
  & \flag{1D-FIN}
  & Extended surfaces / cooling pins ($\varepsilon_f \ge 2$ required to justify fins) \idxformula{fin efficiency!$\eta_f=\tanh(mL_c)/(mL_c)$} \\
$\dfrac{k_e}{\sigma_e T} = k_e \dfrac{\rho_e}{T} = L_0 = \dfrac{\pi^2}{3}\!\left(\dfrac{k_B}{e}\right)^{\!2}$
  & $k_e$: electronic thermal conductivity $(\si{\watt\per\meter\per\kelvin})$; $\sigma_e = 1/\rho_e$: electrical conductivity $(\si{\siemens\per\meter})$; $L_0 = \SI{2.443e-8}{\watt\ohm\per\square\kelvin}$
  & \flag{METALS}
  & Wiedemann--Franz law separating electronic $k_e$ from phonon $k_{\mathrm{ph}} = k - k_e$ \src{S25} \idxformula{Wiedemann-Franz law!$k_e/(\sigma_e T)=L_0$} \\
\end{ftab}
\end{formulacard}

% ----------------------------------------------------------------------------
% L2: COMPACT EXPLANATION & MICROSCOPIC MECHANISMS
% ----------------------------------------------------------------------------
\section{Fourier's Law, Thermal Conductivity Mechanisms \& Diffusivity}
\label{sec:10.1}

\subsection{Macroscopic Law and Sign Convention}
Heat conduction is the diffusive transfer of internal thermal energy from more energetic particles to adjacent less energetic particles without bulk macroscopic fluid motion \src{S12,S13}. In 1822, Jean-Baptiste Joseph Fourier established that the local heat flux $q''$ $[\si{\watt\per\square\meter}]$ is directly proportional to the negative temperature gradient. For steady one-dimensional conduction across a homogeneous plane slab of thickness $L$ $[\si{\meter}]$ and cross-sectional area $A$ $[\si{\square\meter}]$ held at uniform boundary temperatures $T_1$ (at $x=0$) and $T_2$ (at $x=L$):
\begin{equation}
  \dot{Q}_{\mathrm{cond}} = A q'' = -k A \frac{T_2 - T_1}{L - 0} = \frac{k A (T_1 - T_2)}{L} = \frac{\Delta T}{R_{\mathrm{slab}}}, \qquad T(x) = T_1 - \frac{T_1 - T_2}{L}\,x
  \label{eq:10:fourier-slab}
\end{equation}
where $k$ $[\si{\watt\per\meter\per\kelvin}]$ is the material's \textbf{thermal conductivity}. Because heat flows spontaneously from higher temperature to lower temperature ($T_1 > T_2 \implies \dot{Q}_{\mathrm{cond}} > 0$ in the $+x$ direction), the temperature profile $T(x)$ inside a homogeneous plane wall with constant $k$ and no internal heat generation ($\dot{q}_{\mathrm{gen}} = 0$) is strictly \textbf{linear}. If $k(T)$ varies linearly with temperature as $k(T) = k_0[1 + \beta_k(T - T_0)]$, the exact steady conduction rate is obtained simply by evaluating $k$ at the arithmetic mean temperature $T_m = (T_1 + T_2)/2$: $\dot{Q} = k(T_m) A (T_1 - T_2)/L$ \src{S12}.

\subsection{Microscopic Transport Carriers and the Wiedemann--Franz Law}
\idxtopic{Wiedemann-Franz law}
\idxtopic{Lorenz number}
In any condensed or gaseous medium, total thermal conductivity is the sum of contributions from mobile charge carriers (free electrons $k_e$) and quantized lattice vibrations or molecular collisions (phonons/molecules $k_{\mathrm{ph}}$):
\begin{equation}
  k_{\mathrm{tot}} = k_e + k_{\mathrm{ph}}, \qquad k_{\mathrm{kinetic}} = \frac{1}{3} C_V v_{\mathrm{carrier}} \lambda_{\mathrm{mfp}}
  \label{eq:10:carrier-sum}
\end{equation}
where $C_V$ $[\si{\joule\per\cubic\meter\per\kelvin}]$ is the carrier volumetric heat capacity, $v_{\mathrm{carrier}}$ is the carrier speed (Fermi velocity $v_F \sim \SI{e6}{\meter\per\second}$ for electrons; sound speed $c_s \sim \SI{5e3}{\meter\per\second}$ for phonons; mean thermal speed $\bar{v} = \sqrt{8RT/(\pi M)}$ for gases), and $\lambda_{\mathrm{mfp}}$ is the carrier mean free path \src{S12,S15}.

\begin{enumerate}
  \item \textbf{Pure Metals \& Alloys ($k \approx \SIrange{10}{427}{\watt\per\meter\per\kelvin}$):} Conduction is dominated by free conduction electrons ($k_e \gg k_{\mathrm{ph}}$). In 1853, Gustav Wiedemann and Rudolph Franz discovered that at a given absolute temperature $T$, the ratio of electronic thermal conductivity $k_e$ to electrical conductivity $\sigma_e = 1/\rho_e$ is nearly universal across metals. Sommerfeld's quantum free-electron theory yields the exact \textbf{Sommerfeld Lorenz number} $L_0$ (tested on the 2026--2027 Cornell Invitational \src{S25}):
  \begin{equation}
    \frac{k_e}{\sigma_e T} = k_e \frac{\rho_e}{T} = L_0 = \frac{\pi^2}{3}\!\left(\frac{k_B}{e}\right)^{\!2} = \SI{2.44300e-8}{\watt\ohm\per\square\kelvin}
    \label{eq:10:wiedemann-franz}
  \end{equation}
  Alloying a pure metal (e.g., adding $\SI{4}{\percent}\ \mathrm{Cu}$ to pure $\mathrm{Al}$ to make Duralumin, or $\mathrm{Cr/Ni}$ to iron to make AISI 304 stainless steel) introduces impurity scattering centers that sharply reduce $\lambda_{\mathrm{mfp}}$ for electrons, dropping $k$ from $\SI{237}{\watt\per\meter\per\kelvin}$ (pure $\mathrm{Al}$) to $\SI{164}{\watt\per\meter\per\kelvin}$ (Duralumin) and from $\SI{80}{\watt\per\meter\per\kelvin}$ (pure $\mathrm{Fe}$) to $\SI{13.8}{\watt\per\meter\per\kelvin}$ (AISI 304 stainless steel) \src{S12}.
  \item \textbf{Crystalline vs.\ Amorphous Nonmetals ($k \approx \SIrange{0.02}{2000}{\watt\per\meter\per\kelvin}$):} In stiff, low-atomic-mass crystals with few defects, long phonon mean free paths produce extraordinary thermal conductivity without free electrons: Type IIb diamond has $k = \SI{1350}{\watt\per\meter\per\kelvin}$ (and pyrolytic graphite $\parallel$ layer planes has $k = \SI{2000}{\watt\per\meter\per\kelvin}$ at $\SI{20}{\celsius}$, exceeding silver's $\SI{427}{\watt\per\meter\per\kelvin}$ by $>4.6\times$!), silicon has $k = \SI{153}{\watt\per\meter\per\kelvin}$, and sapphire ($\mathrm{Al_2O_3}$) has $k = \SI{46}{\watt\per\meter\per\kelvin}$ \src{S12}. Conversely, in amorphous glasses ($k_{\mathrm{Pyrex}} = \SI{1.3}{\watt\per\meter\per\kelvin}$, $k_{\mathrm{soda\text{-}lime}} = \SI{1.1}{\watt\per\meter\per\kelvin}$) and polymers ($k = \SIrange{0.15}{0.35}{\watt\per\meter\per\kelvin}$), structural disorder restricts $\lambda_{\mathrm{mfp}}$ to a single interatomic bond spacing ($\sim \SI{0.5}{\nano\meter}$).
  \item \textbf{Liquids ($k \approx \SIrange{0.10}{0.68}{\watt\per\meter\per\kelvin}$ for nonmetallic liquids) and Gases ($k \approx \SIrange{0.006}{0.18}{\watt\per\meter\per\kelvin}$):} Liquid water ($k = \SI{0.598}{\watt\per\meter\per\kelvin}$ at $\SI{20}{\celsius}$, rising to $\SI{0.631}{\watt\per\meter\per\kelvin}$ at $\SI{40}{\celsius}$ \src{S09,S12}) is the most conductive common nonmetallic liquid because of its low molar mass and hydrogen-bonded network, yet it is still $\sim 650\times$ less conductive than copper. In ideal gases, kinetic theory gives $k_{\mathrm{gas}} \approx \frac{1}{3}\rho c_v \bar{v}\lambda_{\mathrm{mfp}} \propto \frac{\sqrt{T/M}}{d_{\mathrm{mol}}^2}$, so gas conductivity \textbf{increases with temperature} ($\propto \sqrt{T}$), is \textbf{independent of pressure} down to the Knudsen vacuum regime ($\lambda_{\mathrm{mfp}} \sim L$), and is highest for light gases ($\mathrm{H_2}$: $\SI{0.182}{\watt\per\meter\per\kelvin}$; $\mathrm{He}$: $\SI{0.151}{\watt\per\meter\per\kelvin}$ vs.\ stagnant air: $\SI{0.02623}{\watt\per\meter\per\kelvin}$ and $\mathrm{CO_2}$: $\SI{0.0166}{\watt\per\meter\per\kelvin}$ at $\SI{300}{\kelvin}$ \src{S10,S12}).
\end{enumerate}

\begin{table}[htbp]
\centering
\caption{Verified Thermal Conductivity $k$, Density $\rho$, Specific Heat $c_p$, Diffusivity $\alpha$, and Effusivity $b$ at $\SI{20}{\celsius}$--$\SI{27}{\celsius}$ (Lienhard Tables A.1, A.2, A.3, A.6 \src{S12}).}
\label{tab:10:k-hierarchy}
\footnotesize
\setlength{\tabcolsep}{3.0pt}
\renewcommand{\arraystretch}{1.12}
\begin{tabularx}{\linewidth}{@{}L{0.27\linewidth} r r r r r X@{}}
\toprule
\textbf{Material (\SI{20}{\celsius} unless noted)} & \textbf{$\rho$ $(\si{\kilogram\per\cubic\meter})$} & \textbf{$c_p$ $(\si{\joule\per\kilogram\per\kelvin})$} & \textbf{$k$ $(\si{\watt\per\meter\per\kelvin})$} & \textbf{$\alpha$ $(\si{\square\meter\per\second})$} & \textbf{$b = \sqrt{k\rho c_p}$} & \textbf{Carrier / 2027 Role} \\
\midrule
Pyrolytic graphite ($\parallel$ planes, \SI{27}{\celsius}) & $2200$ & $710$ & $2000$ & $1.28\times 10^{-3}$ & $55\,890$ & Anisotropic phonons ($\perp$: $k=\SI{9.5}{\watt\per\meter\per\kelvin}$) \\
Diamond (Type IIb) & $3250$ & $510$ & $1350$ & $8.14\times 10^{-4}$ & $47\,300$ & Sp$^3$ covalent phonon crystal \\
Silver ($99.99+\%$ pure) & $10\,524$ & $236$ & $427$ & $1.719\times 10^{-4}$ & $32\,570$ & Highest-$k$ pure metal \\
Copper (pure) & $8954$ & $384$ & $398$ & $1.157\times 10^{-4}$ & $37\,000$ & Thin non-contact absorber foil \\
Gold (pure) & $19\,320$ & $129$ & $318$ & $1.276\times 10^{-4}$ & $28\,150$ & Non-oxidizing IR reflector \\
Aluminum (pure / household foil) & $2707$ & $905$ & $237$ & $9.61\times 10^{-5}$ & $24\,100$ & Primary 2027 reflector ($\varepsilon \approx 0.035$) \\
Aluminum alloy 6061-T6 & $2700$ & $896$ & $167$ & $6.90\times 10^{-5}$ & $20\,100$ & Structural flashing/sheet \\
Brass ($70\%\ \mathrm{Cu}, 30\%\ \mathrm{Zn}$) & $8522$ & $385$ & $109$ & $3.32\times 10^{-5}$ & $18\,910$ & Calorimetry cylinder standard \\
Iron (pure) / Carbon steel ($0.5\%\ \mathrm{C}$) & $7833$ & $465$ & $54.0$ & $1.47\times 10^{-5}$ & $14\,030$ & Fastener thermal-bridge hazard! \\
Lead (pure) & $11\,373$ & $130$ & $35.0$ & $2.34\times 10^{-5}$ & $7193$ & Low-$c_p$ calorimetry weight \\
Stainless steel (AISI 304) & $8000$ & $400$ & $13.8$ & $4.31\times 10^{-6}$ & $6645$ & Low-$k$ probe sheath stem \\
Ice Ih (\SI{0}{\celsius}) & $917$ & $2100$ & $2.215$ & $1.15\times 10^{-6}$ & $2065$ & $3.7\times$ higher $k$ than liquid water! \\
Borosilicate glass (Pyrex, \SI{60}{\celsius}) & $2210$ & $753$ & $1.30$ & $7.81\times 10^{-7}$ & $1471$ & Selective greenhouse glazing \\
Liquid water (\SI{27}{\celsius} / \SI{300}{\kelvin}) & $996.6$ & $4179$ & $0.613$ & $1.47\times 10^{-7}$ & $1598$ & Target working fluid ($\mathrm{Bi}_z > 0.1$) \\
Polyethylene (LDPE film / HDPE) & $920$ & $2100$ & $0.330$ & $1.71\times 10^{-7}$ & $798$ & Legal evaporation cling wrap \\
Polycarbonate (Lexan, \SI{23}{\celsius}) & $1200$ & $1250$ & $0.290$ & $1.93\times 10^{-7}$ & $660$ & Rigid glazing window \\
Polypropylene (\texttt{\#5 PP} clear cup) & $905$ & $1900$ & $0.185$ & $1.08\times 10^{-7}$ & $564$ & ES 9 oz party cup wall \\
Cardboard (solid sheet / fluted) & $790 / 180$ & $1340$ & $0.14 / 0.055$ & $1.32\times 10^{-7}$ & $385 / 115$ & Legal structural box \& baffles \\
Plywood (Douglas fir, \SI{24}{\celsius}) & $550$ & $1200$ & $0.120$ & $1.82\times 10^{-7}$ & $281$ & Outer framing baseplate \\
Balsa wood ($\perp$ grain, \SI{15}{\celsius}) & $100$ & $1600$ & $0.050$ & $3.13\times 10^{-7}$ & $89.4$ & Knife-edge 3-point cup support \\
Sheep wool / Cotton batting & $145 / 45$ & $1500 / 1340$ & $0.050 / 0.039$ & $2.30\times 10^{-7}$ & $104 / 48.5$ & Legal fibrous cavity fill \\
Corkboard (medium $\rho$, \SI{30}{\celsius}) & $170$ & $1850$ & $0.040$ & $1.27\times 10^{-7}$ & $112$ & High-strength base insulator \\
Stagnant air (\SI{27}{\celsius}, $1\,\mathrm{atm}$) & $1.176$ & $1007$ & $0.02623$ & $2.213\times 10^{-5}$ & $5.57$ & Optimal $10\text{--}12\text{ mm}$ foil air gap \\
\bottomrule
\end{tabularx}
\end{table}

% ----------------------------------------------------------------------------
% L3: DERIVATIONS, COMPOSITE NETWORKS, CONTACT RESISTANCE, R_CRIT & FINS
% ----------------------------------------------------------------------------
\section{One-Dimensional Geometries, Composite Networks \& Thermal Bridges}
\label{sec:10.2}

\begin{derivation}[Algebra-Only (\texttt{5.d.iv} Compliant) Derivation of Cylindrical and Spherical Shell Resistances \src{S06,S12}]
For radial steady-state conduction without internal heat generation across a spherical shell ($A(r) = 4\pi r^2$) or cylindrical shell ($A(r) = 2\pi r \ell$), energy conservation requires the total radial heat transfer rate $\dot{Q}$ $[\si{\watt}]$ crossing any radius $r \in [r_i, r_o]$ to be independent of $r$.
\begin{enumerate}
  \item \textbf{Spherical Shell via Discrete Concentric Sub-Shells:} Subdivide $[r_i, r_o]$ into $N$ thin concentric shells with radii $r_0 = r_i < r_1 < \dots < r_N = r_o$. For a thin shell $m$ of thickness $\Delta r_m = r_m - r_{m-1}$ and geometric-mean surface area $A_{\mathrm{gm},m} = 4\pi r_{m-1}r_m$, the plane-slab resistance is $\Delta R_m = \frac{r_m - r_{m-1}}{4\pi k r_{m-1}r_m} = \frac{1}{4\pi k}\!\left(\frac{1}{r_{m-1}} - \frac{1}{r_m}\right)$. Summing the $N$ series shells telescopes algebraically without calculus:
  \begin{equation}
    R_{\mathrm{sph}} = \sum_{m=1}^N \Delta R_m = \frac{1}{4\pi k}\sum_{m=1}^N \left(\frac{1}{r_{m-1}} - \frac{1}{r_m}\right) = \frac{1}{4\pi k}\!\left(\frac{1}{r_i} - \frac{1}{r_o}\right) = \frac{r_o - r_i}{4\pi k r_i r_o} = \frac{L}{k A_{\mathrm{gm}}}
    \label{eq:10:rsph}
  \end{equation}
  where $L = r_o - r_i$ and $A_{\mathrm{gm}} = \sqrt{A_i A_o} = 4\pi r_i r_o$ is the exact \textbf{geometric mean} of inner and outer surface areas.
  \item \textbf{Cylindrical Shell via Log-Mean Area ($A_{\mathrm{lm}}$):} Subdivide $[r_i, r_o]$ into $N$ geometric ratio shells with $r_m / r_{m-1} = (r_o/r_i)^{1/N} = 1 + \epsilon$. Each thin shell has $\Delta r_m / r_m \approx \epsilon = \frac{1}{N}\ln(r_o/r_i)$, so summing $N$ series shells yields Lienhard Eq.~(2.22) \src{S12}:
  \begin{equation}
    R_{\mathrm{cyl}} = \frac{\ln(r_o / r_i)}{2\pi k \ell} = \frac{r_o - r_i}{k A_{\mathrm{lm}}}, \qquad A_{\mathrm{lm}} = \frac{A_o - A_i}{\ln(A_o / A_i)} = \frac{2\pi\ell(r_o - r_i)}{\ln(r_o / r_i)}
    \label{eq:10:rcyl}
  \end{equation}
  When the wall is thin ($\delta = r_o - r_i \ll r_i$), the algebraic series $\ln(1 + \delta/r_i) \approx \frac{\delta}{r_i} - \frac{\delta^2}{2r_i^2}$ reduces Eq.~\eref{eq:10:rcyl} directly to the flat-slab formula $R_{\mathrm{cyl}} \approx \frac{\delta}{k (2\pi r_m \ell)}$ with arithmetic mean radius $r_m = (r_i + r_o)/2$ (error $<1\%$ for $r_o/r_i < 1.2$).
\end{enumerate}
\end{derivation}

\subsection{Series--Parallel Composite Walls, Wiener Bounds and Fastener Thermal Bridges}
\idxtopic{composite wall!series and parallel}
\idxtopic{thermal bridge!fasteners and studs}
When heat flows through a composite structure containing both series layers and parallel paths (such as wood studs or metal screws penetrating an insulation wall, shown in \fref{fig:10:composite-wall}):
\begin{itemize}
  \item \textbf{Series Layers (1-D Isothermal Planes $\perp$ Heat Flow):} Every layer carries the same heat flow $\dot{Q}$. For $n$ plane layers of thicknesses $L_i$ and conductivities $k_i$ exposed to fluid temperatures $T_{\infty,\mathrm{in}}$ (with $h_{\mathrm{in}}$) and $T_{\infty,\mathrm{out}}$ (with $h_{\mathrm{out}}$):
  \begin{equation}
    R_{\mathrm{tot,ser}} = \frac{1}{h_{\mathrm{in}}A} + \sum_{i=1}^n \frac{L_i}{k_i A} + \sum R_{c,i} + \frac{1}{h_{\mathrm{out}}A}, \qquad k_{\mathrm{eff,ser}} = \frac{\sum_{i=1}^n L_i}{\sum_{i=1}^n (L_i / k_i)} \quad \text{(Weighted Harmonic Mean)}
    \label{eq:10:series-keff}
  \end{equation}
  The interface temperature $T_m$ after resistance $R_m$ is found by the voltage-divider ratio:
  \begin{equation}
    T_m = T_{\infty,\mathrm{in}} - \dot{Q}\!\left(R_{\mathrm{in}} + \sum_{i=1}^m R_i\right) = T_{\infty,\mathrm{in}} - \left(T_{\infty,\mathrm{in}} - T_{\infty,\mathrm{out}}\right)\frac{R_{\mathrm{in}} + \sum_{i=1}^m R_i}{R_{\mathrm{tot,ser}}}
    \label{eq:10:node-temp}
  \end{equation}
  \item \textbf{Parallel Paths / Thermal Bridges (Heat Flow $\parallel$ Layer Interfaces):} Parallel paths of cross-sectional areas $A_j$ and thickness $L$ share the same temperature drop $\Delta T$, so thermal conductances $G_j = 1/R_j = k_j A_j / L$ add directly:
  \begin{equation}
    \frac{1}{R_{\mathrm{tot,par}}} = \sum_{j=1}^m \frac{1}{R_j} = \frac{1}{L}\sum_{j=1}^m k_j A_j, \qquad k_{\mathrm{eff,par}} = \frac{\sum_{j=1}^m k_j A_j}{\sum_{j=1}^m A_j} \quad \text{(Area-Weighted Arithmetic Mean)}
    \label{eq:10:parallel-keff}
  \end{equation}
  By the Cauchy--Schwarz inequality, for any heterogeneous composite with volume fractions $f_j = A_j/A_{\mathrm{tot}}$, the true 3-D effective conductivity is bounded by the **Wiener bounds**: $k_{\mathrm{eff,ser}} \le k_{\mathrm{eff,3D}} \le k_{\mathrm{eff,par}}$. Because $k_{\mathrm{steel}} / k_{\mathrm{cork}} = 54.0 / 0.040 = 1350$, a steel screw occupying just $\SI{0.17}{\percent}$ of a wall's area conducts \textbf{$2.26\times$ more heat than the remaining $\SI{99.83}{\percent}$ of cork insulation combined}!
\end{itemize}

% ---------------- FIGURE 28: fig:10:composite-wall ---------------------------
\begin{center}
\begin{tikzpicture}[font=\scriptsize, >=Latex, scale=0.95]
  % Left panel: Physical cross-section of series-parallel composite wall
  \node[font=\sffamily\bfseries\scriptsize, color=navy] at (2.3, 4.35) {(a) Physical Composite Wall with Fastener Bridge};
  % Layer 1: Cork (x=0 to 1.1)
  \fill[gray!22] (0,0.4) rectangle (1.1,3.8);
  \draw[thick] (0,0.4) rectangle (1.1,3.8);
  \node[rotate=90, align=center] at (0.55, 2.75) {Layer 1: Cork\\[-1pt]$L_1, k_1$};
  % Layer 2: Parallel Air/Batt (top) + Thermal Bridge Stud/Bolt (bottom)
  \fill[gray!08] (1.1,1.35) rectangle (2.9,3.8);
  \draw[thick] (1.1,1.35) rectangle (2.9,3.8);
  \node[align=center] at (2.0, 2.65) {Layer 2a: Insulation\\($A_{2a}, L_2, k_{2a}$ low)};
  \fill[gray!55] (1.1,0.4) rectangle (2.9,1.35);
  \draw[thick] (1.1,0.4) rectangle (2.9,1.35);
  \node[align=center, font=\sffamily\bfseries\scriptsize] at (2.0, 0.88) {2b: Bolt/Stud Bridge\\[-1pt]($A_{2b}, L_2, k_{2b} \gg k_{2a}$)};
  % Layer 3: Cardboard shell (x=2.9 to 3.9)
  \fill[gray!30] (2.9,0.4) rectangle (3.9,3.8);
  \draw[thick] (2.9,0.4) rectangle (3.9,3.8);
  \node[rotate=90, align=center] at (3.4, 2.1) {Layer 3: Shell\\[-1pt]$L_3, k_3$};
  % Contact interface dashed line at x=1.1
  \draw[rust, very thick, densely dashed] (1.1,0.4) -- (1.1,3.8);
  \node[rust, font=\sffamily\bfseries\scriptsize, anchor=north] at (1.1, 0.38) {$R_c = \frac{1}{h_c A}$};
  % Temperature profile curve overlaid
  \draw[navy, ultra thick] (-0.7,3.55) .. controls (-0.25,3.55) and (-0.1,3.35) .. (0,3.25)
    -- (1.1,2.75) -- (1.1,2.45) -- (2.9,1.30) -- (3.9,0.95)
    .. controls (4.05,0.85) and (4.25,0.70) .. (4.65,0.70);
  \fill[navy] (-0.7,3.55) circle (1.6pt) node[above right, xshift=-4pt] {$T_{\infty,\mathrm{in}}, h_{\mathrm{in}}$};
  \fill[navy] (0,3.25) circle (1.5pt) node[above left] {$T_{s,1}$};
  \fill[rust] (1.1,2.75) circle (1.4pt);
  \fill[rust] (1.1,2.45) circle (1.4pt) node[right, font=\scriptsize] {$\Delta T_c$};
  \fill[navy] (3.9,0.95) circle (1.5pt) node[above right] {$T_{s,3}$};
  \fill[navy] (4.65,0.70) circle (1.6pt) node[above left, xshift=4pt] {$T_{\infty,\mathrm{out}}, h_{\mathrm{out}}$};
  \draw[->, very thick, rust] (-0.55, 1.95) -- (0.35, 1.95) node[midway, above] {$\dot{Q}_{\mathrm{tot}}$};

  % Right panel: Equivalent Thermal-Electrical Resistance Network
  \node[font=\sffamily\bfseries\scriptsize, color=navy] at (10.3, 4.35) {(b) Equivalent Series--Parallel Thermal Resistance Circuit};
  \coordinate (Nin) at (5.5, 2.2);
  \coordinate (Ns1) at (7.0, 2.2);
  \coordinate (N1a) at (8.5, 2.2);
  \coordinate (N1b) at (9.9, 2.2);
  \coordinate (N2)  at (12.1, 2.2);
  \coordinate (Ns3) at (13.6, 2.2);
  \coordinate (Nout) at (15.0, 2.2);

  \draw[thick] (Nin) -- (5.75,2.2);
  \draw[thick] (5.75,2.0) rectangle (6.75,2.4);
  \node at (6.25,2.2) {$R_{\mathrm{in}}$};
  \draw[thick] (6.75,2.2) -- (Ns1);

  \draw[thick] (Ns1) -- (7.25,2.2);
  \draw[thick] (7.25,2.0) rectangle (8.25,2.4);
  \node at (7.75,2.2) {$R_1=\frac{L_1}{k_1 A}$};
  \draw[thick] (8.25,2.2) -- (N1a);

  \draw[thick, rust] (N1a) -- (8.75,2.2);
  \draw[thick, rust] (8.75,2.0) rectangle (9.65,2.4);
  \node[rust] at (9.20,2.2) {$R_c=\frac{1}{h_c A}$};
  \draw[thick, rust] (9.65,2.2) -- (N1b);
  % Parallel split
  \draw[thick] (N1b) -- (9.9, 3.15) -- (10.3, 3.15);
  \draw[thick] (10.3, 3.15) rectangle (11.7, 3.55);
  \node at (11.0, 3.35) {$R_{2a} = \frac{L_2}{k_{2a}A_{2a}}$};
  \draw[thick] (11.7, 3.15) -- (12.1, 3.15) -- (N2);

  \draw[thick] (N1b) -- (9.9, 1.25) -- (10.3, 1.25);
  \fill[gray!35] (10.3, 0.85) rectangle (11.7, 1.25);
  \draw[thick] (10.3, 0.85) rectangle (11.7, 1.25);
  \node[font=\bfseries\scriptsize] at (11.0, 1.05) {$R_{2b} = \frac{L_2}{k_{2b}A_{2b}}$};
  \draw[thick] (11.7, 1.25) -- (12.1, 1.25) -- (N2);
  \node[rust, font=\sffamily\bfseries\scriptsize] at (11.0, 0.52) {$\dot{Q}_{\mathrm{bridge}} = \frac{\Delta T_2}{R_{2b}} \gg \dot{Q}_{\mathrm{ins}}$};

  \draw[thick] (N2) -- (12.4, 2.2);
  \draw[thick] (12.4, 2.0) rectangle (13.3, 2.4);
  \node at (12.85, 2.2) {$R_3$};
  \draw[thick] (13.3, 2.2) -- (Ns3);
  \draw[thick] (Ns3) -- (13.9, 2.2);
  \draw[thick] (13.9, 2.0) rectangle (14.7, 2.4);
  \node at (14.3, 2.2) {$R_{\mathrm{out}}$};
  \draw[thick] (14.7, 2.2) -- (Nout);

  % Nodes
  \foreach \pt/\lbl/\pos in {Nin/$T_{\infty,\mathrm{in}}$/above, Ns1/$T_{s,1}$/below, N1a/$T_{1^-,c}$/below, N1b/$T_{1^+,c}$/above, N2/$T_{23}$/above, Ns3/$T_{s,3}$/below, Nout/$T_{\infty,\mathrm{out}}$/above} {
    \fill[navy] (\pt) circle (1.8pt) node[\pos, yshift=1pt] {\lbl};
  }
  \node[align=center, draw=navy, fill=boxbg, rounded corners=1.5pt, inner sep=3pt] at (10.3, -0.25)
    {$\dot{Q}_{\mathrm{tot}} = \dfrac{T_{\infty,\mathrm{in}} - T_{\infty,\mathrm{out}}}{R_{\mathrm{in}} + R_1 + R_c + \left(R_{2a}^{-1} + R_{2b}^{-1}\right)^{-1} + R_3 + R_{\mathrm{out}}} = U A \left(T_{\infty,\mathrm{in}} - T_{\infty,\mathrm{out}}\right)$};
\end{tikzpicture}
\captionof{figure}{Series--parallel composite wall with interfacial contact resistance $R_c = 1/(h_c A)$ and a high-conductivity fastener thermal bridge $R_{2b}$ in parallel with insulation $R_{2a}$, paired with its exact thermal resistance circuit \src{S12}.}
\label{fig:10:composite-wall}
\idxfig{composite wall!series-parallel thermal bridge network}
\end{center}

\section{Interfacial Contact Resistance \& Critical Radius of Insulation}
\label{sec:10.3}
\idxtopic{contact resistance!interfacial conductance $h_c$}
\idxtopic{critical radius of insulation}

\subsection{Interfacial Contact Conductance ($h_c$)}
Even optically smooth solid surfaces touch only at discrete microscopic asperities ($1\text{--}2\%$ of nominal area $A$ under moderate pressure); the remaining $98\text{--}99\%$ gap is filled with trapped air ($k_{\mathrm{air}} = \SI{0.0262}{\watt\per\meter\per\kelvin}$) or vacuum. Heat crosses the interface by parallel solid-spot constriction conduction and interstitial fluid conduction, producing a discrete temperature jump $\Delta T_c = T_{1^-} - T_{1^+} = \dot{Q} R_c = \dot{Q}/(h_c A)$ \src{S12}. \tref{tab:10:contact-hc} compiles verified contact conductances $h_c$ from Lienhard Table~2.1 (p.~76 \src{S12}):

\begin{table}[htbp]
\centering
\caption{Verified Interfacial Contact Conductances $h_c$ and Unit-Area Contact Resistances $R_c'' = 1/h_c$ for Normal Surface Finishes at Moderate Contact Pressure ($1\text{--}10\ \text{atm}$ unless noted; Lienhard Table 2.1, p.~76 \src{S12}).}
\label{tab:10:contact-hc}
\footnotesize
\setlength{\tabcolsep}{5pt}
\renewcommand{\arraystretch}{1.12}
\begin{tabularx}{\linewidth}{@{}L{0.42\linewidth} r r X@{}}
\toprule
\textbf{Contacting Interface Pair \& Interstitial Condition} & \textbf{$h_c$ $(\si{\watt\per\square\meter\per\kelvin})$} & \textbf{$R_c'' = 1/h_c$ $(\si{\square\meter\kelvin\per\watt})$} & \textbf{Physical \& Device Implication} \\
\midrule
Iron / aluminum ($70\ \text{atm}$ high contact pressure, air) & $45\,000$ & $2.22\times 10^{-5}$ & Plastic deformation flattens asperities at high $P$ \\
Copper / copper ($1\text{--}10\ \text{atm}$, air interstices) & $10\,000\text{--}25\,000$ & $(0.40\text{--}1.00)\times 10^{-4}$ & Soft ductile metal yields low contact resistance \\
Aluminum / aluminum ($1\text{--}10\ \text{atm}$, air interstices) & $2200\text{--}12\,000$ & $(0.83\text{--}4.55)\times 10^{-4}$ & Native $\mathrm{Al_2O_3}$ film adds constriction resistance \\
Graphite / metals ($1\text{--}10\ \text{atm}$, air interstices) & $3000\text{--}6000$ & $(1.67\text{--}3.33)\times 10^{-4}$ & Conformable dry thermal interface pad \\
Stainless steel / stainless steel ($1\text{--}10\ \text{atm}$, air) & $2000\text{--}3700$ & $(2.70\text{--}5.00)\times 10^{-4}$ & Hard asperities; negligible only if slab $L \gg \SI{3}{\milli\meter}$ \\
Ceramic / metals ($1\text{--}10\ \text{atm}$, air interstices) & $1500\text{--}8500$ & $(1.18\text{--}6.67)\times 10^{-4}$ & Ceramic heater / metal heat-sink interface \\
Ceramic / ceramic ($1\text{--}10\ \text{atm}$, air interstices) & $500\text{--}3000$ & $(0.33\text{--}2.00)\times 10^{-3}$ & Brittle point contacts with wide air pockets \\
Stainless steel / stainless steel (evacuated interstices) & $200\text{--}1100$ & $(0.91\text{--}5.00)\times 10^{-3}$ & Removing interstitial gas drops $h_c$ by $3\text{--}10\times$! \\
Aluminum / aluminum (low pressure, evacuated interstices) & $100\text{--}400$ & $(0.25\text{--}1.00)\times 10^{-2}$ & Vacuum multi-layer insulation (MLI) contact blocker \\
Unclamped plastic cup on knife-edge balsa/cork cradle ($\sim 0.01\ \text{atm}$) & $15\text{--}50$ & $(2.0\text{--}6.7)\times 10^{-2}$ & 2027 cup cradle: tiny contact $A_c < \SI{0.25}{\square\centi\meter}$ blocks leak! \\
\bottomrule
\end{tabularx}
\end{table}

\subsection{Critical Radius of Insulation ($r_{\mathrm{crit}}$) — Theory and 2027 9 oz Cup Immunity}
When a cylinder of outer solid radius $r_i$ (at temperature $T_i$) is wrapped with an insulation layer of conductivity $k_{\mathrm{ins}}$ up to outer radius $r_o$ exposed to ambient air at $T_\infty$ with combined convection + radiation coefficient $h_{\mathrm{ext}} = h_{\mathrm{conv}} + h_{\mathrm{rad}}$, the total thermal resistance is:
\begin{equation}
  R_{\mathrm{tot,cyl}}(r_o) = R_{\mathrm{cond}}(r_o) + R_{\mathrm{conv}}(r_o) = \frac{\ln(r_o / r_i)}{2\pi k_{\mathrm{ins}} \ell} + \frac{1}{2\pi r_o \ell h_{\mathrm{ext}}}
  \label{eq:10:rcrit-sum}
\end{equation}
Increasing $r_o$ \textbf{increases} $R_{\mathrm{cond}}$ logarithmically, but \textbf{decreases} $R_{\mathrm{conv}}$ as $1/r_o$ because the outer heat-loss surface area $2\pi r_o \ell$ expands! Setting the rate of change of $R_{\mathrm{tot,cyl}}(r_o)$ with respect to $r_o$ to zero ($\frac{1}{r_o k_{\mathrm{ins}}} - \frac{1}{r_o^2 h_{\mathrm{ext}}} = 0$, Lienhard Example~2.6 \& Eq.~2.27, p.~70 \src{S12}) gives the exact \textbf{critical radius of insulation} where $R_{\mathrm{tot}}$ is a \textbf{minimum} and heat loss $\dot{Q}$ is a \textbf{maximum}:
\begin{equation}
  \boxed{r_{\mathrm{crit,cyl}} = \frac{k_{\mathrm{ins}}}{h_{\mathrm{ext}}} \quad (\mathrm{Bi}_{r_{\mathrm{crit}}} = 1), \qquad r_{\mathrm{crit,sph}} = \frac{2k_{\mathrm{ins}}}{h_{\mathrm{ext}}} \quad (\mathrm{Bi}_{r_{\mathrm{crit}}} = 2)}
  \label{eq:10:rcrit-boxed}
\end{equation}
\begin{itemize}
  \item \textbf{Case 1 ($r_i < r_{\mathrm{crit}}$ — Thin Wire / Capillary Paradox):} Adding insulation initially \textbf{increases} heat dissipation until $r_o = r_{\mathrm{crit}}$ (useful for cooling current-carrying electrical wires!), and does not reduce heat loss below the bare-tube value until $r_o > r_{\mathrm{cross}}$ where $R_{\mathrm{tot}}(r_{\mathrm{cross}}) = R_{\mathrm{tot}}(r_i)$ (plotted in \fref{fig:10:cylindrical-rcrit}).
  \item \textbf{Case 2 ($r_i > r_{\mathrm{crit}}$ — 2027 9 oz Cup \& Standard Pipes):} For the ES-provided 9 oz party cup ($r_i \approx \SIrange{30}{37.5}{\milli\meter}$) surrounded by compliant insulation ($k_{\mathrm{ins}} \approx \SI{0.040}{\watt\per\meter\per\kelvin}$) in room air ($h_{\mathrm{ext}} \approx \SI{10}{\watt\per\square\meter\per\kelvin}$), $r_{\mathrm{crit}} = \frac{0.040}{10} = \SI{0.0040}{\meter} = \SI{4.0}{\milli\meter}$. Because $r_i = \SI{37.5}{\milli\meter} = 9.4\,r_{\mathrm{crit}} \gg r_{\mathrm{crit}}$, we are far to the right of the minimum in \fref{fig:10:cylindrical-rcrit}: \textbf{every millimeter of insulation added around the 9 oz cup monotonically increases $R_{\mathrm{tot}}$ and reduces heat loss!}
\end{itemize}

% ---------------- FIGURE 29: fig:10:cylindrical-rcrit ------------------------
\begin{center}
\begin{tikzpicture}[font=\scriptsize]
  \begin{axis}[
    width=0.86\linewidth, height=5.0cm,
    xlabel={Normalized Outer Insulation Radius $r_o / r_i$ (Lienhard Example 2.6: $r_i = \SI{2.5}{\milli\meter},\; k = \SI{0.074}{\watt\per\meter\per\kelvin},\; h = \SI{20}{\watt\per\square\meter\per\kelvin}$)},
    ylabel={Thermal Resistance $R_t$ $(\si{\kelvin\per\watt\per\meter})$},
    xmin=1.0, xmax=4.5, ymin=0, ymax=4.5,
    xtick={1.0, 1.48, 2.0, 2.323, 3.0, 3.5, 4.0, 4.5},
    xticklabels={$1.0$, $1.48\ (r_{\mathrm{crit}})$, $2.0$, $2.32\ (r_{\mathrm{cross}})$, $3.0$, $3.5$, $4.0$, $4.5$},
    ytick={0, 1.0, 2.0, 3.0, 4.0},
    grid=both, major grid style={gray!30}, minor grid style={gray!15},
    legend pos=north west, legend style={font=\scriptsize, fill=white, draw=navy},
    clip=true
  ]
    % R_cond = ln(x) / (2*pi*0.074) = 2.15076 * ln(x)
    \addplot[teal, thick, densely dashed, domain=1.0:4.5, samples=36] {2.15076 * ln(x)};
    \addlegendentry{$R_{\mathrm{cond}}(r_o) = \dfrac{\ln(r_o/r_i)}{2\pi k \ell}$}

    % R_conv = 1 / (2*pi*20*0.0025*x) = 3.18310 / x
    \addplot[amber, thick, densely dotted, domain=1.0:4.5, samples=36] {3.18310 / x};
    \addlegendentry{$R_{\mathrm{conv}}(r_o) = \dfrac{1}{2\pi r_o \ell h}$}

    % R_tot = R_cond + R_conv
    \addplot[navy, ultra thick, domain=1.0:4.5, samples=36] {2.15076 * ln(x) + 3.18310 / x};
    \addlegendentry{$R_{\mathrm{tot}}(r_o) = R_{\mathrm{cond}} + R_{\mathrm{conv}}$}

    % Horizontal bare-tube baseline R_tot(1) = 3.1831
    \draw[rust, thick, dashed] (axis cs:1.0,3.1831) -- (axis cs:2.3232,3.1831);
    \draw[rust, thick, dotted] (axis cs:1.48,0) -- (axis cs:1.48,2.9936);
    \draw[rust, thick, dotted] (axis cs:2.3232,0) -- (axis cs:2.3232,3.1831);
    \fill[rust] (axis cs:1.0,3.1831) circle (2pt);
    \fill[rust] (axis cs:1.48,2.9936) circle (2.2pt) node[below right, text=rust, font=\bfseries\scriptsize] {Min $R_{\mathrm{tot}}=\SI{2.99}{\kelvin\per\watt}$ ($\dot{Q}_{\max}$ at $\mathrm{Bi}=1$)};
    \fill[navy] (axis cs:2.3232,3.1831) circle (2pt) node[above right, text=navy, font=\bfseries\scriptsize] {Crossover $r_o/r_i = 2.323$};
    \node[align=left, fill=boxbg, draw=slate, rounded corners=1pt, inner sep=2.5pt] at (axis cs:3.45,1.25)
      {\textbf{2027 9 oz Cup Regime:}\\$r_{\mathrm{cup}}=\SI{37.5}{\milli\meter} \gg r_{\mathrm{crit}}=\SI{4.0}{\milli\meter}$\\($\mathrm{Bi}_{\mathrm{cup}} \approx 9.4 > 1 \implies$ insulation\\always increases $R_{\mathrm{tot}}$!)};
  \end{axis}
\end{tikzpicture}
\captionof{figure}{Total cylindrical thermal resistance $R_{\mathrm{tot}}(r_o) = R_{\mathrm{cond}} + R_{\mathrm{conv}}$ versus radius ratio $r_o/r_i$ for Lienhard Example 2.6 \src{S12}. When $r_i < r_{\mathrm{crit}} = k/h$, resistance passes through a minimum at $r_o = r_{\mathrm{crit}}$ ($\mathrm{Bi}=1$) before crossing the bare-pipe resistance at $r_o/r_i = 2.323$.}
\label{fig:10:cylindrical-rcrit}
\idxfig{critical radius of insulation!resistance minimum curve}
\end{center}

\section{Extended Surfaces (Fins) \& Experimental Conductivity Determination}
\label{sec:10.4}
\idxtopic{fins and extended surfaces!efficiency and effectiveness}
\idxtopic{thermal conductivity!experimental measurement}

\subsection{Algebraic Fin Design Equations (When Convection Limits Heat Transfer)}
When $R_{\mathrm{conv}} = 1/(hA)$ is much larger than internal solid conduction resistance, extended surfaces (\textbf{fins}) of cross-sectional area $A_c$, perimeter $P$, length $L$, and conductivity $k$ are attached to a base at excess temperature $\theta_b = T_b - T_\infty$ (Lienhard §4.5, pp.~177--196 \src{S12}). Define the \textbf{fin parameter} $m$ $[\si{\per\meter}]$ and corrected tip length $L_c = L + A_c/P$ ($L_c = L + t/2$ for a rectangular plate fin of thickness $t$; $L_c = L + D/4$ for a cylindrical pin fin of diameter $D$):
\begin{equation}
  m = \sqrt{\frac{h P}{k A_c}}, \qquad m L_c = \sqrt{\mathrm{Bi}_{\mathrm{fin}}}\,\frac{L_c}{\sqrt{A_c/P}}, \qquad \dot{Q}_{\mathrm{fin}} = \sqrt{h P k A_c}\,\theta_b \tanh(m L_c) = \eta_f h A_f \theta_b
  \label{eq:10:fin-q}
\end{equation}
\begin{itemize}
  \item \textbf{Fin Efficiency ($\eta_f$) and Overall Finned-Array Efficiency ($\eta_o$):}
  \begin{equation}
    \eta_f = \frac{\dot{Q}_{\mathrm{fin}}}{h A_f \theta_b} = \frac{\tanh(m L_c)}{m L_c} \in (0, 1], \qquad \eta_o = 1 - \frac{A_{\mathrm{fins}}}{A_{\mathrm{total}}}\left(1 - \eta_f\right), \qquad R_{\mathrm{th,array}} = \frac{1}{\eta_o h A_{\mathrm{total}}}
    \label{eq:10:fin-eff}
  \end{equation}
  \item \textbf{Fin Effectiveness ($\varepsilon_f$):} Ratio of fin heat transfer rate to the heat transfer rate from the bare footprint area $A_c$ without the fin:
  \begin{equation}
    \varepsilon_f = \frac{\dot{Q}_{\mathrm{fin}}}{h A_c \theta_b} = \sqrt{\frac{k P}{h A_c}}\,\tanh(m L_c) \xrightarrow{m L_c \gtrsim 2.65\ (\text{long fin})} \varepsilon_{f,\infty} = \sqrt{\frac{k P}{h A_c}}
    \label{eq:10:fin-effectiveness}
  \end{equation}
  \textbf{Engineering Rules:} (1) Practical fins require $\varepsilon_f \ge 2$ ($k P \ge 4 h A_c$); (2) place fins on the \textbf{low-$h$ gas side}, never on a boiling/condensing side where $h$ is huge; (3) choose thin, closely spaced high-$k$ copper/aluminum fins to maximize perimeter-to-area ratio $P/A_c = 2/t$; (4) beyond $mL_c \approx 2.65$, $\tanh(mL_c) > 0.99$, so making a fin longer adds mass without adding heat transfer!
\end{itemize}

\subsection{Experimental Determination of Thermal Conductivity ($k$)}
\begin{enumerate}
  \item \textbf{Symmetrical Guarded Hot-Plate Apparatus (ASTM C177 — Standard for Low-$k$ Insulators):} An electrical main heater of area $A$ dissipating measured DC power $P_{\mathrm{elec}} = V I$ is sandwiched symmetrically between two identical flat specimen slabs of thickness $L$, backed by water-cooled cold plates at $T_c$. A coplanar concentric \textbf{guard ring heater} around the main heater's perimeter is servo-controlled to the exact main-heater temperature $T_h$ ($T_{\mathrm{guard}} - T_h = 0$), forcing zero radial edge heat leak! Because half of $P_{\mathrm{elec}}$ flows through each identical specimen slab:
  \begin{equation}
    \dot{Q}_{\mathrm{one\,slab}} = \frac{P_{\mathrm{elec}}}{2} = \frac{V I}{2} \implies k = \frac{(V I / 2)\,L}{A\,(T_h - T_c)} = \frac{V I L}{2 A \Delta T}
    \label{eq:10:guarded-hot-plate}
  \end{equation}
  \item \textbf{Comparative Cut-Bar Method (Axial Series Stack):} An unknown disk of thickness $L_x$ and conductivity $k_x$ is clamped in series between two calibrated reference cylinders of known conductivity $k_{\mathrm{ref}}$ and thickness $L_{\mathrm{ref}}$ inside an insulated guard shell. By series heat-flux continuity ($q''_{\mathrm{ref}} = q''_x$):
  \begin{equation}
    k_{\mathrm{ref}}\frac{\Delta T_{\mathrm{ref}}}{L_{\mathrm{ref}}} = k_x \frac{\Delta T_x}{L_x} \implies k_x = k_{\mathrm{ref}}\,\frac{L_x}{L_{\mathrm{ref}}}\,\frac{\Delta T_{\mathrm{ref}}}{\Delta T_x}
    \label{eq:10:comparative-bar}
  \end{equation}
  \item \textbf{Searle's Bar Method (High-$k$ Metal Rod with Flowing Water Calorimeter):} One end of an insulated metal bar (area $A$, thermocouple spacing $L$, temperatures $T_1, T_2$) is steam-heated while the cold end is cooled by steady water flow $\dot{m}_w$ warming from $T_{w,\mathrm{in}}$ to $T_{w,\mathrm{out}}$:
  \begin{equation}
    \dot{Q} = \dot{m}_w c_{p,w}\left(T_{w,\mathrm{out}} - T_{w,\mathrm{in}}\right) = k A \frac{T_1 - T_2}{L} \implies k = \frac{\dot{m}_w c_{p,w}(T_{w,\mathrm{out}} - T_{w,\mathrm{in}}) L}{A(T_1 - T_2)}
    \label{eq:10:searles-bar}
  \end{equation}
\end{enumerate}

\subsection{Multi-Dimensional Conduction Shape Factors ($S$) \& Temperature-Dependent Conductivity $k(T)$}
\idxtopic{conduction shape factor!3D edges, corners, buried bodies}
\idxtopic{thermal conductivity!temperature dependence $k(T)$}
When steady conduction occurs through two- or three-dimensional geometries (such as a buried spherical/cylindrical tank or the 12 edges and 8 corners of a thick-walled cubic collector enclosure), the total heat flow between isothermal surfaces $T_1$ and $T_2$ is expressed via the \textbf{conduction shape factor} $S$ $[\si{\meter}]$ (Lienhard §5.7 Table~5.4, p.~242 \src{S12}):
\begin{equation}
  \dot{Q} = S k \left(T_1 - T_2\right), \qquad R_{\mathrm{th,3D}} = \frac{1}{S k} \quad [\si{\kelvin\per\watt}]
  \label{eq:10:shape-factor}
\end{equation}
\begin{itemize}
  \item \textbf{Langmuir / Lienhard 3-D Thick-Walled Box Shape Factor (Walls + 12 Edges + 8 Corners):} For a closed rectangular box of uniform wall thickness $L_w \le a_{\mathrm{in}}/5$ with inner surface area $A_{\mathrm{in}}$ and 12 inner edge lengths $\ell_{\mathrm{edge},k}$:
  \begin{equation}
    S_{\mathrm{box,3D}} = \underbrace{\frac{A_{\mathrm{in}}}{L_w}}_{\text{6 plane walls}} + \underbrace{0.54\sum_{k=1}^{12} \ell_{\mathrm{edge},k}}_{\text{12 edges }(>\SI{13}{\percent}\text{ boost!})} + \underbrace{0.15\times 8\,L_w}_{\text{8 corners}}
    \label{eq:10:langmuir-box}
  \end{equation}
  \item \textbf{Canonical Buried \& Concentric Shape Factors (Lienhard Table 5.4, p.~242 \src{S12}):}
  \begin{enumerate}
    \item \emph{Isothermal sphere of radius $r$ buried at depth $z$ below an isothermal ground plane:} $S = \dfrac{4\pi r}{1 - r/(2z)} \xrightarrow{z \to \infty} 4\pi r$ ($\mathrm{Nu}_D \to 2$).
    \item \emph{Horizontal cylinder of radius $r$ and length $\ell \gg r$ buried at depth $z$:} $S = \dfrac{2\pi\ell}{\operatorname{arccosh}(z/r)} \approx \dfrac{2\pi\ell}{\ln(2z/r)}$ (for $z > 3r$).
    \item \emph{Concentric spheres of radii $r_i, r_o$:} $S = \dfrac{4\pi r_i r_o}{r_o - r_i}$; \emph{Concentric cylinders of length $\ell$:} $S = \dfrac{2\pi\ell}{\ln(r_o/r_i)}$.
  \end{enumerate}
  \item \textbf{Temperature-Dependent Thermal Conductivity $k(T)$ (Kirchhoff Transform \& Mean-Conductivity Theorem):} When $k(T)$ varies with temperature across $[T_2, T_1]$, separating variables in $\dot{Q} = -k(T)A\,\frac{dT}{dx}$ shows that $\dot{Q}$ equals the exact conduction rate computed with the \textbf{integral-average thermal conductivity} $\bar{k}$:
  \begin{equation}
    \dot{Q} = \frac{A}{L}\int_{T_2}^{T_1} k(T)\,dT = \bar{k}\,\frac{A(T_1 - T_2)}{L}, \qquad \text{For linear } k(T) = k_0\left[1 + \beta_k(T - T_0)\right]:\quad \bar{k} = k\!\left(\frac{T_1 + T_2}{2}\right)
    \label{eq:10:kt-average}
  \end{equation}
\end{itemize}

\begin{table}[htbp]
\centering
\caption{Temperature Dependence of Metallic Thermal Conductivities $k(T)$ $[\si{\watt\per\meter\per\kelvin}]$ from $\SI{-170}{\celsius}$ to $\SI{800}{\celsius}$ (Lienhard Table A.1, pp.~739--740 \src{S12}).}
\label{tab:10:k-vs-temp}
\footnotesize
\setlength{\tabcolsep}{4.5pt}
\renewcommand{\arraystretch}{1.10}
\begin{tabularx}{\linewidth}{@{}L{0.26\linewidth} r r r r r r r r X@{}}
\toprule
\textbf{Metal / Alloy} & \textbf{$\SI{-170}{\celsius}$} & \textbf{$\SI{-100}{\celsius}$} & \textbf{$\SI{0}{\celsius}$} & \textbf{$\SI{20}{\celsius}$} & \textbf{$\SI{100}{\celsius}$} & \textbf{$\SI{200}{\celsius}$} & \textbf{$\SI{400}{\celsius}$} & \textbf{$\SI{600}{\celsius}$} & \textbf{Trend $dk/dT$ \& Mechanism} \\
\midrule
Aluminum (pure) & $302$ & $242$ & $236$ & $237$ & $240$ & $238$ & $228$ & $215$ & Flat plateau $\SIrange{0}{200}{\celsius}$ \\
Aluminum 6061-T6 & --- & --- & $166$ & $167$ & $172$ & $177$ & $180^*$ & --- & Alloy: $k$ rises slightly with $T$ ($^*$\SI{300}{\celsius}) \\
Copper (pure) & $483$ & $420$ & $401$ & $398$ & $391$ & $389$ & $378$ & $366$ & Decreases slowly ($\beta_k \approx \SI{-2.3e-4}{\per\kelvin}$) \\
Brass ($70\mathrm{Cu}$--$30\mathrm{Zn}$) & $73$ & $89$ & $106$ & $109$ & $133$ & $143$ & $147$ & --- & Strong positive slope ($+24\%$ by $\SI{100}{\celsius}$) \\
Iron (pure) & $132$ & $98$ & $84$ & $80$ & $72$ & $63$ & $50$ & $39$ & Phonon+magnon scattering drops $k$ \\
Carbon steel ($0.5\%\ \mathrm{C}$) & --- & --- & $55$ & $54$ & $52$ & $48$ & $42$ & $35$ & Monotonic decrease with $T$ \\
Stainless AISI 304 & --- & --- & $13.8$ & $13.8$ & $15.0$ & $17.2$ & $21.0$ & $25.0$ & Austenitic alloy: $k$ increases $+81\%$! \\
Silver ($99.99+\%$ pure) & $449$ & $431$ & $428$ & $427$ & $422$ & $417$ & $401$ & $386$ & Highest metallic $k$ across all $T$ \\
Tungsten (filament) & $235$ & $223$ & $182$ & $178$ & $166$ & $153$ & $134$ & $125$ & Drops to $\sim\SI{105}{\watt\per\meter\per\kelvin}$ at $\SI{2400}{\kelvin}$ \\
\bottomrule
\end{tabularx}
\end{table}

% ----------------------------------------------------------------------------
% L4: WORKED EXAMPLES (7 Fully Verified Examples Matching check_sec10_13.py)
% ----------------------------------------------------------------------------
\section{Worked Competition Examples (§10)}
\label{sec:10.5}

\begin{workedexample}[Multi-Layer Collector Wall \& Interface Temperature Profiling]{2027 Device Cradle / Regional--National}{ex:10.1}
A flat side panel of a 2027 heat-collector enclosure has cross-sectional area $A = \SI{0.0400}{\square\meter}$ and separates warm internal cavity air at $T_{\infty,\mathrm{in}} = \SI{46.0}{\celsius}$ ($h_{\mathrm{in}} = \SI{12.0}{\watt\per\square\meter\per\kelvin}$) from room air at $T_{\infty,\mathrm{out}} = \SI{21.0}{\celsius}$ ($h_{\mathrm{out}} = \SI{8.00}{\watt\per\square\meter\per\kelvin}$). From inside to outside, the wall consists of: (1) a corkboard liner ($L_1 = \SI{4.00}{\milli\meter},\; k_1 = \SI{0.0400}{\watt\per\meter\per\kelvin}$), (2) a stagnant foil-bounded dead-air gap ($L_2 = \SI{10.0}{\milli\meter},\; k_{\mathrm{eff},2} = \SI{0.0275}{\watt\per\meter\per\kelvin}$), and (3) an outer fluted corrugated cardboard shell ($L_3 = \SI{4.00}{\milli\meter},\; k_3 = \SI{0.0550}{\watt\per\meter\per\kelvin}$). Find (a) the total thermal resistance $R_{\mathrm{tot}}$ and overall heat transfer coefficient $U$, (b) the steady heat loss $\dot{Q}$, and (c) the temperature $T_{12}$ at the cork--air gap interface.
\tcblower
\textbf{Solution.} All 5 thermal resistances act in series across $A = \SI{0.0400}{\square\meter}$:
\begin{align*}
  R_{\mathrm{in}} &= \frac{1}{h_{\mathrm{in}}A} = \frac{1}{(12.0)(0.0400)} = \SI{2.0833}{\kelvin\per\watt}, &
  R_1 &= \frac{L_1}{k_1 A} = \frac{0.00400}{(0.0400)(0.0400)} = \SI{2.5000}{\kelvin\per\watt} \\
  R_2 &= \frac{L_2}{k_{\mathrm{eff},2}A} = \frac{0.0100}{(0.0275)(0.0400)} = \SI{9.0909}{\kelvin\per\watt}, &
  R_3 &= \frac{L_3}{k_3 A} = \frac{0.00400}{(0.0550)(0.0400)} = \SI{1.8182}{\kelvin\per\watt} \\
  R_{\mathrm{out}} &= \frac{1}{h_{\mathrm{out}}A} = \frac{1}{(8.00)(0.0400)} = \SI{3.1250}{\kelvin\per\watt}
\end{align*}
(a) Summing in series: $R_{\mathrm{tot}} = 2.0833 + 2.5000 + 9.0909 + 1.8182 + 3.1250 = \mathbf{\SI{18.617}{\kelvin\per\watt}}$, and $U = \frac{1}{R_{\mathrm{tot}}A} = \frac{1}{(18.6174)(0.0400)} = \mathbf{\SI{1.343}{\watt\per\square\meter\per\kelvin}}$ ($R_{\mathrm{val,tot}} = 1/U = \SI{0.7447}{\square\meter\kelvin\per\watt}$).\\
(b) Steady heat transfer rate: $\dot{Q} = \dfrac{T_{\infty,\mathrm{in}} - T_{\infty,\mathrm{out}}}{R_{\mathrm{tot}}} = \dfrac{46.0 - 21.0}{18.6174} = \mathbf{\SI{1.343}{\watt}}$.\\
(c) By Eq.~\eref{eq:10:node-temp}, the cork--air gap interface temperature is:
\[
  T_{12} = T_{\infty,\mathrm{in}} - \dot{Q}\left(R_{\mathrm{in}} + R_1\right) = 46.0 - (1.34283)(2.0833 + 2.5000) = 46.0 - 6.1546 = \mathbf{\SI{39.85}{\celsius}}
\]
\verifiedby{Independent $5\times 5$ nodal temperature matrix solve in \texttt{check\_sec10\_13.py} (\texttt{EX-10.1-Q}, \texttt{EX-10.1-T12}): $\dot{Q} = \SI{1.34283}{\watt}$, $T_{12} = \SI{39.8454}{\celsius}$ ($\text{rel\_err} = 0.00$).}
\end{workedexample}

\begin{workedexample}[Fastener Thermal Bridge Short-Circuit in an Insulated Box]{Parallel Resistance Trap}{ex:10.2}
A cork insulation wall of total area $A_{\mathrm{tot}} = \SI{0.0300}{\square\meter}$ and thickness $L = \SI{15.0}{\milli\meter}$ ($k_{\mathrm{cork}} = \SI{0.0400}{\watt\per\meter\per\kelvin}$) is fastened using four carbon-steel M4 bolts (each of diameter $d = \SI{4.00}{\milli\meter}$, $k_{\mathrm{steel}} = \SI{54.0}{\watt\per\meter\per\kelvin}$) passing completely through the $\SI{15.0}{\milli\meter}$ thickness. Compare the wall's pure-cork thermal resistance $R_{\mathrm{pure}}$ against its actual parallel thermal resistance $R_{\mathrm{eq}}$ with the 4 bolts, and find the percentage of heat that leaks through the bolts.
\tcblower
\textbf{Solution.} Pure cork resistance: $R_{\mathrm{pure}} = \dfrac{L}{k_{\mathrm{cork}}A_{\mathrm{tot}}} = \dfrac{0.0150}{(0.0400)(0.0300)} = \mathbf{\SI{12.50}{\kelvin\per\watt}}$.\\
With 4 bolts, total bolt cross-sectional area is $A_{\mathrm{bolts}} = 4\left(\frac{\pi}{4}d^2\right) = \pi(\SI{0.00400}{\meter})^2 = \SI{5.0265e-5}{\square\meter}$ (just $\SI{0.1676}{\percent}$ of $A_{\mathrm{tot}}$!), leaving $A_{\mathrm{cork}} = 0.0300 - 0.00005027 = \SI{0.029950}{\square\meter}$. The parallel conductances $G = 1/R = kA/L$ are:
\[
  G_{\mathrm{cork}} = \frac{(0.0400)(0.029950)}{0.0150} = \SI{0.07987}{\watt\per\kelvin}, \qquad
  G_{\mathrm{bolts}} = \frac{(54.0)(\SI{5.0265e-5}{})}{0.0150} = \SI{0.18096}{\watt\per\kelvin}
\]
\[
  R_{\mathrm{eq}} = \frac{1}{G_{\mathrm{cork}} + G_{\mathrm{bolts}}} = \frac{1}{0.07987 + 0.18096} = \frac{1}{0.26082} = \mathbf{\SI{3.834}{\kelvin\per\watt}}, \qquad
  f_{\mathrm{bolts}} = \frac{G_{\mathrm{bolts}}}{G_{\mathrm{tot}}} = \frac{0.18096}{0.26082} = \mathbf{\SI{69.38}{\percent}}
\]
Four tiny steel screws occupying $\SI{0.17}{\percent}$ of the wall area destroy $\SI{69.3}{\percent}$ of the wall's thermal resistance! Replacing steel bolts with balsa dowels ($k = \SI{0.050}{\watt\per\meter\per\kelvin}$) or glue keeps $R_{\mathrm{eq}} = \SI{12.49}{\kelvin\per\watt}$.
\verifiedby{Parallel conductance sum in \texttt{check\_sec10\_13.py} (\texttt{EX-10.2-R-STEEL}, \texttt{EX-10.2-FRAC}): $R_{\mathrm{eq}} = \SI{3.83404}{\kelvin\per\watt}$, $f_{\mathrm{bolts}} = 0.69379$.}
\end{workedexample}

\begin{workedexample}[Interfacial Contact Resistance Between Pressed Stainless Steel Plates]{Lienhard Ex.\ 2.9 Extension \src{S12}}{ex:10.3}
Two AISI 304 stainless steel plates, each of thickness $L = \SI{15.0}{\milli\meter}$, area $A = \SI{0.0200}{\square\meter}$, and conductivity $k = \SI{15.0}{\watt\per\meter\per\kelvin}$, are pressed together in air with interfacial contact conductance $h_c = \SI{2500}{\watt\per\square\meter\per\kelvin}$ (\tref{tab:10:contact-hc}). An overall temperature drop $\Delta T_{\mathrm{tot}} = \SI{40.0}{\kelvin}$ is maintained across the outer faces. Find the steady heat transfer rate $\dot{Q}$ and the interfacial temperature jump $\Delta T_c$.
\tcblower
\textbf{Solution.} Each stainless plate has bulk conductive resistance $R_{\mathrm{plate}} = \frac{L}{kA} = \frac{0.0150}{(15.0)(0.0200)} = \SI{0.0500}{\kelvin\per\watt}$. The contact interface adds series resistance $R_c = \frac{1}{h_c A} = \frac{1}{(2500)(0.0200)} = \SI{0.0200}{\kelvin\per\watt}$:
\[
  R_{\mathrm{tot}} = 2 R_{\mathrm{plate}} + R_c = 2(0.0500) + 0.0200 = \SI{0.1200}{\kelvin\per\watt} \implies
  \dot{Q} = \frac{\Delta T_{\mathrm{tot}}}{R_{\mathrm{tot}}} = \frac{40.0}{0.1200} = \mathbf{\SI{333.3}{\watt}}
\]
The temperature jump across the microscopic contact plane is $\Delta T_c = \dot{Q} R_c = (333.33)(0.0200) = \mathbf{\SI{6.67}{\kelvin}}$ ($\SI{16.7}{\percent}$ of $\Delta T_{\mathrm{tot}}$).
\verifiedby{Series resistance check in \texttt{check\_sec10\_13.py} (\texttt{EX-10.3-Q}, \texttt{EX-10.3-DT-INT}): $\dot{Q} = \SI{333.333}{\watt}$, $\Delta T_c = \SI{6.6667}{\kelvin}$.}
\end{workedexample}

\begin{workedexample}[Critical Radius of Insulation \& Crossover Radius vs.\ 9 oz Cup Cradle]{Lienhard Ex.\ 2.6 \src{S12}}{ex:10.4}
A small tube of outer diameter $D_i = \SI{5.00}{\milli\meter}$ ($r_i = \SI{2.50}{\milli\meter}$) is insulated with $85\%$ magnesia ($k = \SI{0.0740}{\watt\per\meter\per\kelvin}$) in ambient air with $h_{\mathrm{ext}} = \SI{20.0}{\watt\per\square\meter\per\kelvin}$ (\fref{fig:10:cylindrical-rcrit}). Find (a) the critical radius $r_{\mathrm{crit}}$ and minimum total resistance per unit length, (b) the crossover outer radius $r_{\mathrm{cross}}$ required before insulation begins reducing heat loss below the bare tube, and (c) compare against the 2027 9 oz cup ($r_{\mathrm{cup}} = \SI{37.5}{\milli\meter},\; k_{\mathrm{cork}} = \SI{0.040}{\watt\per\meter\per\kelvin},\; h_{\mathrm{ext}} = \SI{10.0}{\watt\per\square\meter\per\kelvin}$).
\tcblower
\textbf{Solution.} (a) By Eq.~\eref{eq:10:rcrit-boxed}, $r_{\mathrm{crit}} = \frac{k}{h_{\mathrm{ext}}} = \frac{0.0740}{20.0} = \mathbf{\SI{3.70}{\milli\meter}}$ ($r_{\mathrm{crit}}/r_i = 1.480$). For $\ell = \SI{1.00}{\meter}$, the bare tube resistance is $R_{\mathrm{bare}} = \frac{1}{2\pi r_i \ell h_{\mathrm{ext}}} = \frac{1}{2\pi(0.00250)(1.00)(20.0)} = \mathbf{\SI{3.183}{\kelvin\per\watt}}$, whereas at $r_o = r_{\mathrm{crit}} = \SI{3.70}{\milli\meter}$, $R_{\min} = \frac{\ln(1.480)}{2\pi(0.0740)} + \frac{1}{2\pi(0.00370)(20.0)} = 0.8432 + 2.1508 = \mathbf{\SI{2.994}{\kelvin\per\watt}}$ ($\SI{6.3}{\percent}$ higher heat loss than bare!).\\
(b) Setting $R_{\mathrm{tot}}(r_{\mathrm{cross}}) = R_{\mathrm{bare}} = \SI{3.1831}{\kelvin\per\watt}$ gives $2.15076\ln(r_{\mathrm{cross}}/r_i) + \frac{3.18310}{r_{\mathrm{cross}}/r_i} = 3.18310 \implies \frac{r_{\mathrm{cross}}}{r_i} = \mathbf{2.323}$, so $r_{\mathrm{cross}} = \mathbf{\SI{5.81}{\milli\meter}}$.\\
(c) For the 2027 9 oz cup: $r_{\mathrm{crit,cup}} = \frac{0.040}{10.0} = \SI{4.00}{\milli\meter} \ll r_{\mathrm{cup}} = \SI{37.5}{\milli\meter}$ ($r_{\mathrm{cup}}/r_{\mathrm{crit}} = 9.375$), so adding cork/foil insulation around the cup \textbf{always reduces heat loss}.
\verifiedby{Brent root-finding in \texttt{check\_sec10\_13.py} (\texttt{EX-10.4-RCRIT}): $r_{\mathrm{cross}}/r_i = 2.32321$ ($r_{\mathrm{cross}} = \SI{5.808}{\milli\meter}$).}
\end{workedexample}

\begin{workedexample}[Spherical Shell Conduction \& Geometric-Mean Area $A_{\mathrm{gm}}$]{All Levels}{ex:10.5}
A spherical vessel of inner radius $r_i = \SI{8.00}{\centi\meter}$ and outer radius $r_o = \SI{12.0}{\centi\meter}$ is insulated with cork ($k = \SI{0.0400}{\watt\per\meter\per\kelvin}$). If the inner wall is at $T_i = \SI{80.0}{\celsius}$ and the outer wall is at $T_o = \SI{20.0}{\celsius}$, calculate the steady heat loss $\dot{Q}$ and verify using the geometric-mean area $A_{\mathrm{gm}} = \sqrt{A_i A_o}$.
\tcblower
\textbf{Solution.} Using Eq.~\eref{eq:10:rsph}:
\[
  R_{\mathrm{sph}} = \frac{r_o - r_i}{4\pi k r_i r_o} = \frac{0.120 - 0.080}{4\pi(0.0400)(0.080)(0.120)} = \SI{8.2893}{\kelvin\per\watt} \implies
  \dot{Q} = \frac{80.0 - 20.0}{8.2893} = \mathbf{\SI{7.238}{\watt}}
\]
\verifiedby{Geometric-mean area check in \texttt{check\_sec10\_13.py} (\texttt{EX-10.5-SPH}): $A_i = \SI{0.080425}{\square\meter}$, $A_o = \SI{0.180956}{\square\meter}$, $A_{\mathrm{gm}} = \sqrt{A_i A_o} = \SI{0.120637}{\square\meter} \implies \dot{Q} = k A_{\mathrm{gm}}\Delta T/(r_o - r_i) = \SI{7.23823}{\watt}$.}
\end{workedexample}

\begin{workedexample}[Wiedemann--Franz Law: Electronic vs.\ Alloy Thermal Conductivity]{2027 Cornell Invitational E1 \src{S25}}{ex:10.6}
At $T = \SI{20.0}{\celsius}$ ($\SI{293.15}{\kelvin}$), pure copper has electrical resistivity $\rho_{e,\mathrm{Cu}} = \SI{1.68e-8}{\ohm\meter}$, while structural aluminum alloy 6061-T6 has $\rho_{e,\mathrm{6061}} = \SI{3.99e-8}{\ohm\meter}$. Using the Sommerfeld Lorenz number $L_0 = \frac{\pi^2}{3}(k_B/e)^2 = \SI{2.443e-8}{\watt\ohm\per\square\kelvin}$, compute the theoretical free-electron thermal conductivity $k_e$ of each metal and compare against experimental values ($k_{\mathrm{Cu}} = \SI{398}{\watt\per\meter\per\kelvin}$, $k_{\mathrm{6061}} = \SI{167}{\watt\per\meter\per\kelvin}$, \tref{tab:10:k-hierarchy}).
\tcblower
\textbf{Solution.} Applying Eq.~\eref{eq:10:wiedemann-franz}, $k_e = L_0 T / \rho_e$:
\[
  k_{e,\mathrm{Cu}} = \frac{(\SI{2.4430e-8}{\watt\ohm\per\square\kelvin})(\SI{293.15}{\kelvin})}{\SI{1.68e-8}{\ohm\meter}} = \mathbf{\SI{426.3}{\watt\per\meter\per\kelvin}}, \qquad
  k_{e,\mathrm{6061}} = \frac{(\SI{2.4430e-8}{})(\SI{293.15}{})}{\SI{3.99e-8}{}} = \mathbf{\SI{179.5}{\watt\per\meter\per\kelvin}}
\]
Both predictions agree with experiment within $\sim \SI{7}{\percent}$ (at room temperature, inelastic electron--phonon scattering slightly reduces the effective Lorenz ratio below $L_0$ to $L_{\mathrm{eff}} \approx \SI{2.28e-8}{\watt\ohm\per\square\kelvin}$).
\verifiedby{Exact CODATA evaluation in \texttt{check\_sec10\_13.py} (\texttt{EX-10.6-WF-CU}, \texttt{EX-10.6-WF-6061}): $k_{e,\mathrm{Cu}} = \SI{426.29}{\watt\per\meter\per\kelvin}$, $k_{e,\mathrm{6061}} = \SI{179.49}{\watt\per\meter\per\kelvin}$.}
\end{workedexample}

\begin{workedexample}[Guarded Hot-Plate Experimental Determination of Cork Conductivity]{Lab \& Data Analysis}{ex:10.7}
In a symmetrical ASTM C177 guarded hot-plate test, two identical corkboard disks of thickness $L = \SI{12.0}{\milli\meter}$ and metered main-heater area $A = \SI{0.0250}{\square\meter}$ sandwich a central DC heater powered at $V = \SI{12.00}{\volt}$ and $I = \SI{0.350}{\ampere}$. When the guard ring is balanced ($T_{\mathrm{guard}} = T_h$), steady plate temperatures are $T_h = \SI{54.5}{\celsius}$ and $T_c = \SI{29.3}{\celsius}$. Calculate $k_{\mathrm{cork}}$ and state the error if the student forgets the factor of $1/2$ for bidirectional heat split.
\tcblower
\textbf{Solution.} Total electrical heater power is $P_{\mathrm{elec}} = V I = (12.00)(0.350) = \SI{4.200}{\watt}$. Because heat exits symmetrically through both cork specimens, each slab carries $\dot{Q}_1 = P_{\mathrm{elec}}/2 = \SI{2.100}{\watt}$ across $\Delta T = 54.5 - 29.3 = \SI{25.2}{\kelvin}$:
\[
  k_{\mathrm{cork}} = \frac{\dot{Q}_1 L}{A \Delta T} = \frac{(2.100)(0.0120)}{(0.0250)(25.2)} = \frac{0.0252}{0.630} = \mathbf{\SI{0.0400}{\watt\per\meter\per\kelvin}}
\]
\ifgot{$\SI{0.0800}{\watt\per\meter\per\kelvin}$}{you used full $P_{\mathrm{elec}} = \SI{4.20}{\watt}$ instead of $P_{\mathrm{elec}}/2 = \SI{2.10}{\watt}$ for a two-sided symmetrical guarded hot plate!}
\verifiedby{Checked in \texttt{check\_sec10\_13.py} (\texttt{EX-10.7-GHP}): $k = \SI{0.0400}{\watt\per\meter\per\kelvin}$, matching Lienhard Table A.2 p.~742 \src{S12}.}
\end{workedexample}

\begin{workedexample}[3-D Shape-Factor Conduction Loss of a Cubic Collector Box (Walls + Edges + Corners)]{Lienhard §5.7 \src{S12}}{ex:10.8}
A closed cubic test enclosure of outer side $a_{\mathrm{out}} = \SI{30.0}{\centi\meter}$ ($\SI{0.300}{\meter}$) and uniform wall thickness $L_w = \SI{3.00}{\centi\meter}$ ($\SI{0.0300}{\meter}$, inner side $a_{\mathrm{in}} = a_{\mathrm{out}} - 2L_w = \SI{0.240}{\meter}$) is constructed of corkboard ($k = \SI{0.0400}{\watt\per\meter\per\kelvin}$). When the inner wall surfaces are at $T_1 = \SI{45.0}{\celsius}$ and the outer wall surfaces are at $T_2 = \SI{21.0}{\celsius}$ ($\Delta T = \SI{24.0}{\kelvin}$), calculate (a) the naive 1-D plane-wall heat loss based on the inner surface area $A_{\mathrm{in}} = 6 a_{\mathrm{in}}^2$, and (b) the true 3-D heat loss $\dot{Q}_{\mathrm{3D}}$ including all 12 edges and 8 corners via Eq.~\eref{eq:10:langmuir-box}.
\tcblower
\textbf{Solution.} (a) The 6 inner faces have area $A_{\mathrm{in}} = 6(0.240)^2 = \SI{0.3456}{\square\meter}$, giving $S_{\mathrm{1D}} = \frac{A_{\mathrm{in}}}{L_w} = \frac{0.3456}{0.0300} = \SI{11.520}{\meter}$ and naive 1-D loss $\dot{Q}_{\mathrm{1D}} = S_{\mathrm{1D}}k\Delta T = (11.520)(0.0400)(24.0) = \mathbf{\SI{11.06}{\watt}}$.\\
(b) Including the 12 inner edges (each of length $\ell_{\mathrm{edge}} = \SI{0.240}{\meter}$) and 8 three-dimensional corners via Eq.~\eref{eq:10:langmuir-box}:
\[
  S_{\mathrm{edges}} = 0.54\left(12 \times \SI{0.240}{\meter}\right) = \SI{1.5552}{\meter}, \qquad
  S_{\mathrm{corners}} = 0.15\left(8 \times \SI{0.0300}{\meter}\right) = \SI{0.0360}{\meter}
\]
\[
  S_{\mathrm{box,3D}} = 11.520 + 1.5552 + 0.0360 = \mathbf{\SI{13.1112}{\meter}} \implies
  \dot{Q}_{\mathrm{3D}} = (13.1112)(0.0400)(24.0) = \mathbf{\SI{12.59}{\watt}}
\]
Edge and corner 3-D flaring increases conductive heat leak by $+\SI{13.8}{\percent}$ ($+\SI{1.53}{\watt}$) over the naive inner-area formula!
\verifiedby{Checked in \texttt{check\_sec10\_13.py} (\texttt{EX-10.8-3DBOX}): $S_{\mathrm{box,3D}} = \SI{13.1112}{\meter}$, $\dot{Q}_{\mathrm{3D}} = \SI{12.5868}{\watt}$.}
\end{workedexample}

\begin{workedexample}[Pin-Fin Efficiency $\eta_f$ and Fin Effectiveness $\varepsilon_f$]{Lienhard §4.5 \src{S12}}{ex:10.9}
A cylindrical aluminum alloy 6061-T6 pin fin ($k = \SI{167}{\watt\per\meter\per\kelvin}$) of diameter $D = \SI{4.00}{\milli\meter}$ and protruding length $L = \SI{30.0}{\milli\meter}$ is attached to a base at $T_b = \SI{65.0}{\celsius}$ in forced air at $T_\infty = \SI{25.0}{\celsius}$ ($\theta_b = \SI{40.0}{\kelvin}$) with $h = \SI{25.0}{\watt\per\square\meter\per\kelvin}$. Compute (a) the fin parameter $m$ and corrected length $L_c = L + D/4$, (b) the fin efficiency $\eta_f$, (c) the heat dissipation rate $\dot{Q}_{\mathrm{fin}}$, and (d) the fin effectiveness $\varepsilon_f$.
\tcblower
\textbf{Solution.} Cross-sectional area is $A_c = \frac{\pi}{4}(\SI{0.00400}{\meter})^2 = \SI{1.2566e-5}{\square\meter}$ and perimeter is $P = \pi(\SI{0.00400}{\meter}) = \SI{1.2566e-2}{\meter}$ ($P/A_c = 4/D = \SI{1000}{\per\meter}$).\\
(a) Corrected tip length: $L_c = 30.0 + 4.00/4 = \SI{31.0}{\milli\meter} = \SI{0.0310}{\meter}$. Fin parameter:
\[
  m = \sqrt{\frac{h P}{k A_c}} = \sqrt{\frac{(25.0)(1000)}{167}} = \mathbf{\SI{12.235}{\per\meter}} \implies m L_c = (12.235)(0.0310) = \mathbf{0.3793}
\]
(b) Fin efficiency: $\eta_f = \dfrac{\tanh(0.37929)}{0.37929} = \dfrac{0.36210}{0.37929} = \mathbf{0.9547}$ ($\SI{95.5}{\percent}$).\\
(c) Total effective fin surface area is $A_f = P L_c = (\SI{1.2566e-2}{})(0.0310) = \SI{3.8956e-4}{\square\meter}$, so:
\[
  \dot{Q}_{\mathrm{fin}} = \eta_f h A_f \theta_b = (0.95465)(25.0)(\SI{3.8956e-4}{})(40.0) = \mathbf{\SI{0.3719}{\watt}}
\]
(d) Without the fin, the bare base footprint $A_c$ dissipates $\dot{Q}_{\mathrm{base}} = h A_c \theta_b = (25.0)(\SI{1.2566e-5}{})(40.0) = \SI{0.012566}{\watt}$, giving effectiveness $\varepsilon_f = \dfrac{\dot{Q}_{\mathrm{fin}}}{\dot{Q}_{\mathrm{base}}} = \dfrac{0.3719}{0.012566} = \mathbf{29.59 \gg 2}$.
\verifiedby{Checked in \texttt{check\_sec10\_13.py} (\texttt{EX-10.9-FIN-ETA}, \texttt{EX-10.9-FIN-EPS}): $\eta_f = 0.95465$, $\varepsilon_f = 29.594$.}
\end{workedexample}

\begin{workedexample}[Comparative Cut-Bar Identification of an Unknown Structural Alloy]{Experimental Metrology}{ex:10.10}
In an insulated comparative cut-bar stack, an unknown aluminum alloy cylinder of length $L_x = \SI{30.0}{\milli\meter}$ is placed in axial series with a calibrated $70$--$30$ brass reference cylinder ($k_{\mathrm{ref}} = \SI{109.0}{\watt\per\meter\per\kelvin}$, length $L_{\mathrm{ref}} = \SI{25.0}{\milli\meter}$) of identical diameter. At steady state, thermocouples record axial temperature drops $\Delta T_{\mathrm{ref}} = \SI{18.40}{\kelvin}$ across the brass reference and $\Delta T_x = \SI{14.41}{\kelvin}$ across the unknown specimen. Determine $k_x$ and identify the alloy from \tref{tab:10:k-hierarchy}.
\tcblower
\textbf{Solution.} By Eq.~\eref{eq:10:comparative-bar}, series heat-flux continuity $q''_{\mathrm{ref}} = q''_x$ gives:
\[
  k_x = k_{\mathrm{ref}}\,\frac{L_x}{L_{\mathrm{ref}}}\,\frac{\Delta T_{\mathrm{ref}}}{\Delta T_x} = (109.0)\!\left(\frac{30.0}{25.0}\right)\!\left(\frac{18.40}{14.4115}\right) = \mathbf{\SI{167.0}{\watt\per\meter\per\kelvin}} \longrightarrow \textbf{Aluminum Alloy 6061-T6}
\]
\verifiedby{Checked in \texttt{check\_sec10\_13.py} (\texttt{EX-10.10-CUTBAR}): $k_x = \SI{167.00}{\watt\per\meter\per\kelvin}$.}
\end{workedexample}

% ----------------------------------------------------------------------------
% L5: TRAPS & DIAGNOSTICS + L6: 2027 DEVICE APPLICATION
% ----------------------------------------------------------------------------
\section{Conduction Traps, Diagnostics \& 2027 Collector Cradle Engineering}
\label{sec:10.6}

\begin{trapbox}[Top Conduction \& Thermal Resistance Traps on Science Olympiad Exams]
\idxerr{conduction!adding parallel resistances instead of conductances}
\idxerr{conduction!using whole wall area in parallel branch}
\idxerr{conduction!Imperial vs SI R-value factor 5.678}
\idxerr{conduction!using slab formula on thick cylinder}
\par\smallskip\noindent
\begin{errtab}
\textbf{1. Using total wall area $A_{\mathrm{tot}}$ inside parallel branch $R_j = L/(k_j A_j)$}
  & Gives completely wrong parallel conductance; each parallel path $j$ MUST use its own cross-sectional area $A_j$ ($\sum A_j = A_{\mathrm{tot}}$).
  & Check $\sum f_j = 1$ in $k_{\mathrm{eff,par}} = \sum f_j k_j$. \\
\textbf{2. Adding parallel resistances ($R_1 + R_2$) instead of conductances ($1/R_1 + 1/R_2$)}
  & Predicts that a metal bolt \emph{increases} resistance instead of short-circuiting it!
  & $R_{\mathrm{eq,par}} = \frac{R_1 R_2}{R_1 + R_2}$ is ALWAYS smaller than $\min(R_1, R_2)$. \\
\textbf{3. Confusing US Imperial $R$-value ($\text{ft}^2\cdot^\circ\text{F}\cdot\text{h/BTU}$) with SI $R_{\mathrm{val}}$ ($\si{\square\meter\kelvin\per\watt}$)}
  & Off by a factor of $5.6783\times$ (e.g., fiberglass R-13 in US units is $R_{\mathrm{SI}} = 13 / 5.6783 = \SI{2.29}{\square\meter\kelvin\per\watt}$).
  & Convert first: $R_{\mathrm{SI}} = \SI{0.17611}{} \times R_{\mathrm{US}}$. \\
\textbf{4. Using flat-slab $R = (r_o - r_i)/(2\pi r_i \ell k)$ on a thick cylindrical wall ($r_o/r_i > 1.2$)}
  & Overestimates area (if using $r_o$) or underestimates area (if using $r_i$).
  & Always use $\ln(r_o/r_i)/(2\pi k \ell)$ (log-mean area $A_{\mathrm{lm}}$). \\
\textbf{5. Forgetting the factor of $1/2$ in a symmetrical guarded hot-plate apparatus}
  & Doubles the calculated thermal conductivity $k$.
  & Check whether the heater is sandwiched between two identical slabs ($\dot{Q}_{\mathrm{slab}} = V I / 2$). \\
\textbf{6. Confusing thermal conductivity $k$ with thermal diffusivity $\alpha = k/(\rho c_p)$}
  & Stagnant air has very low $k = \SI{0.0262}{\watt\per\meter\per\kelvin}$, yet its thermal diffusivity $\alpha_{\mathrm{air}} = \SI{2.21e-5}{\square\meter\per\second}$ is \textbf{$150\times$ larger than liquid water's} ($\SI{1.47e-7}{\square\meter\per\second}$) because $\rho_{\mathrm{air}}$ is tiny!
  & Steady flux uses $k$; transient penetration $\delta_t \sim 2\sqrt{\alpha t}$ uses $\alpha$. \\
\end{errtab}
\end{trapbox}

\begin{devicebox}[2027 Heat-Lamp Collector Conduction \& Cup-Cradle Playbook \src{S01,S02,S12}]
\idxdev{cup cradle!3-point knife-edge balsa/cork support}
\idxdev{black paper sleeve trap!low conductivity blocking radiative entry}
\idxdev{Rule 3.e!zero water contact prohibition}
\begin{enumerate}
  \item \textbf{Three-Point Knife-Edge Cup Seat ($\dot{Q}_{\mathrm{base}} < \SI{0.02}{\watt}$):} Never rest the 9 oz cup flat on a solid table or thick wooden floor ($A_{\mathrm{base}} \approx \SI{22}{\square\centi\meter}$ conducts heat downward into the cold floor). Support the cup base on \textbf{three radial balsa or cork knife-edge fins} ($k = \SIrange{0.04}{0.05}{\watt\per\meter\per\kelvin}$, contact footprint $A_c < \SI{0.25}{\square\centi\meter}$, height $L = \SI{35}{\milli\meter}$): $R_{\mathrm{support}} = \frac{0.035}{(0.045)(\SI{2.5e-5}{})} = \SI{31100}{\kelvin\per\watt}$, reducing conductive base heat loss to $<\SI{0.001}{\watt}$!
  \item \textbf{Why Wrapping Black Paper Around the Outside of the Clear Cup Fails:} Matte black construction paper has low thermal conductivity ($k = \SI{0.14}{\watt\per\meter\per\kelvin}$) and poor contact conductance ($h_c \approx \SI{30}{\watt\per\square\meter\per\kelvin}$) against the tapered plastic cup wall. Wrapping black paper around the clear cup \textbf{blocks direct near-IR radiation from passing through the transparent cup wall into the water}, while dissipating absorbed heat outward to the air! Instead, surround the cup with a non-contact specular aluminum-foil cavity ($\varepsilon = 0.04$) or a foil-thin ($\le\SI{35}{\micro\meter}$) high-$k$ copper/aluminum sleeve backed by concentric $\SI{10}{\milli\meter}$ foil-air-foil insulation rings.
  \item \textbf{Rule \texttt{3.e} Zero-Water-Contact Constraint:} Although metallic fins have high fin efficiency $\eta_f$, \textbf{Rule \texttt{3.e} strictly prohibits any part of the device from contacting the water} \src{S01}. All conductive/radiative collection elements must remain outside the cup or in the dry upper air space above the $\SI{100}{\milli\liter}$ waterline ($\ge\SI{15}{\milli\meter}$ clearance).
\end{enumerate}
\end{devicebox}

% ----------------------------------------------------------------------------
% RAPID COMPETITION PROBLEM & SOLUTION BANK (§10)
% ----------------------------------------------------------------------------
\section{Rapid Competition Problem \& Solution Bank (§10)}
\label{sec:10.7}

\begin{problem}{P10.1}{Conduction + Phase Change (Stefan Ice-Layer Growth)}{10.2}{5.c.ii, 5.c.iii}{National / Invitational}{90 s}{Algebra + Energy Balance}{forgetting average thickness $L/2$ or using water $k$ instead of ice $k$}
A calm freshwater lake is at uniform temperature $T_w = \SI{0.00}{\celsius}$ when cold winter air suddenly drops the upper ice surface temperature to a constant $T_{\mathrm{air}} = \SI{-15.0}{\celsius}$ ($\Delta T = \SI{15.0}{\kelvin}$). Ice Ih has thermal conductivity $k_{\mathrm{ice}} = \SI{2.215}{\watt\per\meter\per\kelvin}$ (\tref{tab:10:k-hierarchy}), density $\rho_{\mathrm{ice}} = \SI{917.0}{\kilogram\per\cubic\meter}$, and latent heat of fusion $L_f = \SI{333.5}{\kilo\joule\per\kilogram}$ \src{S11,S12}. Neglecting sensible subcooling of the ice ($c_p\Delta T \ll L_f$), derive the algebraic Stefan law for ice thickness $y(t)$ at time $t$ by equating conductive heat flux across average thickness to latent heat release, and compute the time required to freeze an ice sheet from $y_1 = \SI{2.00}{\centi\meter}$ to $y_2 = \SI{8.00}{\centi\meter}$.
\end{problem}
\begin{solution}{P10.1}
When the ice layer has thickness $y$ and grows by thickness $\Delta y$ over area $A$, the latent heat released at the $y$ interface ($\SI{0}{\celsius}$) is $\Delta Q = \rho_{\mathrm{ice}} A L_f\,\Delta y$, which must be conducted upward through the ice layer to the $\SI{-15.0}{\celsius}$ top surface at rate $\dot{Q} = k_{\mathrm{ice}} A \Delta T / y$. Writing $\rho_{\mathrm{ice}} L_f \frac{\Delta y}{\Delta t} = \frac{k_{\mathrm{ice}}\Delta T}{y_{\mathrm{avg}}}$ where the arithmetic mean thickness between $y_1$ and $y_2$ is $y_{\mathrm{avg}} = \frac{y_1 + y_2}{2}$ gives $y_{\mathrm{avg}}(y_2 - y_1) = \frac{1}{2}(y_2^2 - y_1^2) = \frac{k_{\mathrm{ice}}\Delta T}{\rho_{\mathrm{ice}} L_f}\,\Delta t$:
\[
  \Delta t = \frac{\rho_{\mathrm{ice}} L_f \left(y_2^2 - y_1^2\right)}{2 k_{\mathrm{ice}} \Delta T} = \frac{(917.0)(\SI{3.335e5}{\joule\per\kilogram})\left(0.0800^2 - 0.0200^2\right)}{2(2.215)(15.0)} = \frac{\num{1.8349e6}}{66.45} = \mathbf{\SI{27613}{\second}} = \mathbf{\SI{7.670}{\hour}}
\]
\emph{Quick Check:} Notice $t \propto y^2$: growing from $\SI{0}{\centi\meter}$ to $\SI{8}{\centi\meter}$ takes $16\times$ as long as growing from $\SI{0}{\centi\meter}$ to $\SI{2}{\centi\meter}$, so the increment from $2$ to $\SI{8}{\centi\meter}$ takes $16 - 1 = 15$ units of the first $\SI{2}{\centi\meter}$ freeze time ($\SI{1840.9}{\second} \times 15 = \SI{27613}{\second}$).
\end{solution}

\begin{problem}{P10.2}{Series-Parallel Wiener Bounds \& Anisotropic Laminate}{10.2}{5.c.iii}{State / National}{75 s}{Harmonic vs.\ Arithmetic Mean}{swapping series and parallel conductivity formulas}
A composite heat-spreader laminate is made by alternating $50$ thin sheets of pure copper ($k_{\mathrm{Cu}} = \SI{398}{\watt\per\meter\per\kelvin}$, thickness $t_{\mathrm{Cu}} = \SI{0.200}{\milli\meter}$) with $50$ thin sheets of polyimide film ($k_{\mathrm{PI}} = \SI{0.350}{\watt\per\meter\per\kelvin}$, thickness $t_{\mathrm{PI}} = \SI{0.100}{\milli\meter}$, Lienhard Table~A.2 p.~743 \src{S12}). Assuming perfect bonding ($R_c \to 0$), calculate (a) the in-plane thermal conductivity $k_{\parallel}$ (heat flow parallel to the sheets), (b) the through-thickness thermal conductivity $k_{\perp}$ (heat flow perpendicular to the sheets), and (c) the anisotropy ratio $k_{\parallel} / k_{\perp}$.
\end{problem}
\begin{solution}{P10.2}
Volume fractions are $f_{\mathrm{Cu}} = \frac{0.200}{0.200 + 0.100} = \frac{2}{3}$ and $f_{\mathrm{PI}} = \frac{0.100}{0.300} = \frac{1}{3}$.
\begin{enumerate}
  \item[(a)] \textbf{In-Plane (Parallel Paths, Eq.~\ref{eq:10:parallel-keff}):} $k_{\parallel} = f_{\mathrm{Cu}} k_{\mathrm{Cu}} + f_{\mathrm{PI}} k_{\mathrm{PI}} = \frac{2}{3}(398) + \frac{1}{3}(0.350) = 265.33 + 0.12 = \mathbf{\SI{265.5}{\watt\per\meter\per\kelvin}}$.
  \item[(b)] \textbf{Through-Thickness (Series Layers, Eq.~\ref{eq:10:series-keff}):} $k_{\perp} = \left(\frac{f_{\mathrm{Cu}}}{k_{\mathrm{Cu}}} + \frac{f_{\mathrm{PI}}}{k_{\mathrm{PI}}}\right)^{-1} = \left(\frac{2/3}{398} + \frac{1/3}{0.350}\right)^{-1} = \left(0.001675 + 0.952381\right)^{-1} = \mathbf{\SI{1.048}{\watt\per\meter\per\kelvin}}$.
  \item[(c)] \textbf{Anisotropy Ratio:} $\dfrac{k_{\parallel}}{k_{\perp}} = \dfrac{265.45}{1.0482} = \mathbf{253.3}$. (In-plane heat spreading is dominated by copper; through-plane insulation is dominated by polyimide!).
\end{enumerate}
\end{solution}

\begin{problem}{P10.3}{Spherical vs.\ Cylindrical Critical Radius Paradox}{10.3}{5.c.iii}{Regional / State}{60 s}{Direct Ratio}{forgetting factor of 2 in spherical $r_{\mathrm{crit}} = 2k/h$}
A small spherical thermistor bead of bare diameter $D_i = \SI{2.00}{\milli\meter}$ ($r_i = \SI{1.00}{\milli\meter}$) is coated with an epoxy shell of thermal conductivity $k_{\mathrm{ep}} = \SI{0.220}{\watt\per\meter\per\kelvin}$ (Lienhard Table~A.2 p.~743 \src{S12}) and placed in an airflow with $h_{\mathrm{ext}} = \SI{110}{\watt\per\square\meter\per\kelvin}$. (a) Find the outer epoxy radius $r_{\mathrm{crit,sph}}$ that maximizes heat transfer from the bead. (b) At $r_o = r_{\mathrm{crit,sph}}$, by what percentage does the heat dissipation rate exceed that of the bare $\SI{1.00}{\milli\meter}$-radius bead?
\end{problem}
\begin{solution}{P10.3}
(a) For a sphere (Eq.~\ref{eq:10:rcrit-boxed}), $r_{\mathrm{crit,sph}} = \dfrac{2 k_{\mathrm{ep}}}{h_{\mathrm{ext}}} = \dfrac{2(0.220)}{110} = \SI{0.00400}{\meter} = \mathbf{\SI{4.00}{\milli\meter}}$ ($r_{\mathrm{crit}}/r_i = 4.00$).\\
(b) Bare sphere resistance is $R_{\mathrm{bare}} = \dfrac{1}{4\pi r_i^2 h_{\mathrm{ext}}} = \dfrac{1}{4\pi(0.00100)^2(110)} = \mathbf{\SI{723.43}{\kelvin\per\watt}}$. At $r_o = r_{\mathrm{crit,sph}} = \SI{4.00}{\milli\meter}$:
\[
  R_{\min} = \frac{r_o - r_i}{4\pi k_{\mathrm{ep}} r_i r_o} + \frac{1}{4\pi r_o^2 h_{\mathrm{ext}}} = \frac{0.00300}{4\pi(0.220)(0.00100)(0.00400)} + \frac{723.43}{16} = 271.29 + 45.21 = \mathbf{\SI{316.50}{\kelvin\per\watt}}
\]
Heat transfer increases by $\dfrac{\dot{Q}_{\max}}{\dot{Q}_{\mathrm{bare}}} = \dfrac{R_{\mathrm{bare}}}{R_{\min}} = \dfrac{723.43}{316.50} = 2.2857 \implies \mathbf{+\SI{128.6}{\percent}}$!
\end{solution}

\begin{problem}{P10.4}{Searle's Bar Calorimetric Conductivity Measurement}{10.4}{5.c.iii}{All Levels}{75 s}{Energy Balance}{forgetting radius-to-area $\pi r^2$ or gram-to-kg conversion}
In a Searle's bar experiment, a well-insulated cylindrical brass rod of diameter $D = \SI{3.00}{\centi\meter}$ has two thermocouples spaced $L = \SI{15.0}{\centi\meter}$ apart along its length reading $T_1 = \SI{92.0}{\celsius}$ and $T_2 = \SI{34.0}{\celsius}$. Cooling water flows through the cold-end copper coil at steady mass flow rate $\dot{m}_w = \SI{1.45}{\gram\per\second}$ ($c_{p,w} = \SI{4.184}{\joule\per\gram\per\kelvin}$), entering at $T_{w,\mathrm{in}} = \SI{18.20}{\celsius}$ and exiting at $T_{w,\mathrm{out}} = \SI{22.85}{\celsius}$. Calculate the measured thermal conductivity $k$ of the brass rod.
\end{problem}
\begin{solution}{P10.4}
Cross-sectional area of the rod is $A = \frac{\pi}{4}(\SI{0.0300}{\meter})^2 = \SI{7.0686e-4}{\square\meter}$. Steady heat rate absorbed by the cooling water is:
\[
  \dot{Q} = \dot{m}_w c_{p,w}\left(T_{w,\mathrm{out}} - T_{w,\mathrm{in}}\right) = (\SI{1.45}{\gram\per\second})(\SI{4.184}{\joule\per\gram\per\kelvin})(22.85 - 18.20) = \mathbf{\SI{28.211}{\watt}}
\]
By Eq.~\eref{eq:10:searles-bar}, $k = \dfrac{\dot{Q} L}{A(T_1 - T_2)} = \dfrac{(28.211)(0.150)}{(\SI{7.0686e-4}{})(92.0 - 34.0)} = \dfrac{4.2316}{0.04100} = \mathbf{\SI{103.2}{\watt\per\meter\per\kelvin}}$ (matching $70$--$30$ brass at $\SI{0}{\celsius}$--$\SI{20}{\celsius}$, \tref{tab:10:k-vs-temp}).
\end{solution}

\begin{problem}{P10.5}{Temperature-Dependent Conductivity $k(T) = k_0(1 + \beta_k T)$ in a Refractory Wall}{10.4}{5.c.iii}{National / Cornell E1}{75 s}{Mean-Conductivity Theorem}{evaluating $k$ at only one boundary instead of $T_m$}
An insulating firebrick wall of thickness $L = \SI{0.200}{\meter}$ and area $A = \SI{1.50}{\square\meter}$ has thermal conductivity that varies linearly with temperature as $k(T) = 0.100 + \num{1.25e-4}\,T$ $[\si{\watt\per\meter\per\kelvin}]$, where $T$ is in $^\circ\mathrm{C}$ (Lienhard Table~A.2 p.~741 \src{S12}). If the inner furnace face is held at $T_1 = \SI{900.0}{\celsius}$ and the outer face is at $T_2 = \SI{100.0}{\celsius}$, calculate (a) the steady heat transfer rate $\dot{Q}$, and (b) the exact temperature $T_{\mathrm{mid}}$ at the geometric midplane $x = L/2 = \SI{0.100}{\meter}$.
\end{problem}
\begin{solution}{P10.5}
(a) By Eq.~\eref{eq:10:kt-average}, for linear $k(T)$ the integral average $\bar{k}$ equals $k(T_m)$ at the arithmetic mean temperature $T_m = \frac{1}{2}(900.0 + 100.0) = \SI{500.0}{\celsius}$:
\[
  \bar{k} = 0.100 + (\num{1.25e-4})(500.0) = \mathbf{\SI{0.1625}{\watt\per\meter\per\kelvin}} \implies
  \dot{Q} = \bar{k}\,\frac{A(T_1 - T_2)}{L} = \frac{(0.1625)(1.50)(800.0)}{0.200} = \mathbf{\SI{975.0}{\watt}}
\]
(b) At the midplane $x = L/2$, half of the integrated conductivity potential has elapsed: $\int_{T_{\mathrm{mid}}}^{T_1} k(T)\,dT = \int_{T_2}^{T_{\mathrm{mid}}} k(T)\,dT$. Since $k(T)$ increases with $T$, the hotter half ($x < L/2$) has higher conductivity and therefore a \emph{smaller} temperature drop than the colder half! Solving $0.100(T_{\mathrm{mid}} - 100) + \frac{\num{1.25e-4}}{2}(T_{\mathrm{mid}}^2 - 100^2) = \frac{1}{2}(0.1625)(800) = 65.00$ gives $\num{6.25e-5}T_{\mathrm{mid}}^2 + 0.100 T_{\mathrm{mid}} - 75.625 = 0 \implies T_{\mathrm{mid}} = \mathbf{\SI{541.5}{\celsius}} > \SI{500.0}{\celsius}$.
\end{solution}

\begin{problem}{P10.6}{Coaxial Two-Layer Pipe Insulation \& Optimal Layer Ordering}{10.2}{5.c.iii}{National}{90 s}{Logarithmic Radial Area}{putting low-k layer on the outside instead of the inside}
A hot pipe of outer radius $r_1 = \SI{20.0}{\milli\meter}$ is to be insulated with two concentric $\SI{10.0}{\milli\meter}$-thick cylindrical sleeves (so intermediate radius is $r_2 = \SI{30.0}{\milli\meter}$ and outer radius is $r_3 = \SI{40.0}{\milli\meter}$): Material~A ($k_A = \SI{0.0250}{\watt\per\meter\per\kelvin}$) and Material~B ($k_B = \SI{0.0500}{\watt\per\meter\per\kelvin}$). Should Material~A (the better insulator) be placed on the inner shell ($[r_1, r_2]$) or the outer shell ($[r_2, r_3]$) to maximize total thermal resistance, and by what percentage do the two arrangements differ?
\end{problem}
\begin{solution}{P10.6}
In cylindrical or spherical geometry, heat flux $q''(r) = \dot{Q}/(2\pi r \ell)$ is highest at the smallest radius $r_1$, so the \textbf{lowest-conductivity material ($k_A$) must ALWAYS be placed on the innermost layer ($[r_1, r_2]$)} where cross-sectional area is smallest and $\ln(r_2/r_1) = \ln(1.50) = 0.40547 > \ln(r_3/r_2) = \ln(4/3) = 0.28768$! For $\ell = \SI{1.00}{\meter}$:
\begin{align*}
  R_{\mathrm{A\text{-}inner}} &= \frac{1}{2\pi}\!\left[\frac{\ln(30/20)}{0.0250} + \frac{\ln(40/30)}{0.0500}\right] = \frac{16.2186 + 5.7536}{2\pi} = \mathbf{\SI{3.4970}{\kelvin\per\watt}} \\
  R_{\mathrm{B\text{-}inner}} &= \frac{1}{2\pi}\!\left[\frac{\ln(30/20)}{0.0500} + \frac{\ln(40/30)}{0.0250}\right] = \frac{8.1093 + 11.5073}{2\pi} = \mathbf{\SI{3.1221}{\kelvin\per\watt}}
\end{align*}
Placing the superior insulator on the inside increases conductive resistance by $+\SI{12.01}{\percent}$! (In a flat plane wall, order does not change $R_{\mathrm{tot}}$, but in radial geometry inner placement always wins).
\end{solution}

\par\smallskip\noindent
\begin{minipage}{\linewidth}
\captionof{table}{Conduction Shape Factors $S$ ($\dot{Q} = S k \Delta T$), Internal Volumetric Heating Profiles ($\dot{q}$ $[\si{\watt\per\cubic\meter}]$), and Porous/Dispersed Composite Conductivity Laws (Lienhard §2.3 \& Table~5.4, pp.~62--66, 253 \src{S12}).}
\label{tab:10:shape-factors}
\idxtopic{conduction shape factor!buried cylinder and sphere}
\idxtopic{internal heat generation!parabolic slab, cylinder, sphere}
\idxtopic{Maxwell-Eucken effective conductivity!dispersed spherical pores}
\footnotesize
\setlength{\tabcolsep}{4pt}
\renewcommand{\arraystretch}{1.12}
\begin{tabularx}{\linewidth}{@{}L{0.31\linewidth} L{0.41\linewidth} X@{}}
\toprule
\textbf{Geometry / Physical System} & \textbf{Exact Closed-Form Expression ($S$, $\Delta T_{\max}$, or $k_{\mathrm{eff}}$)} & \textbf{Conditions \& Notes} \\
\midrule
Isothermal sphere (radius $r_1$) buried at depth $z$ below isothermal surface &
  $S = \dfrac{4\pi r_1}{1 - r_1/(2z)} \xrightarrow{z \gg r_1} 4\pi r_1$ &
  Semi-infinite medium ($k$); exact image-source solution \\
Horizontal cylinder (radius $r_1$, length $\ell$) buried at axis depth $z$ &
  $S = \dfrac{2\pi \ell}{\operatorname{arcosh}(z/r_1)} \approx \dfrac{2\pi \ell}{\ln(2z/r_1)}$ \quad ($z > 3r_1$) &
  Requires $\ell \gg r_1$; for $z/r_1 = 5$, $\ln(10) = 2.3026$ \\
Two parallel cylinders (radii $r_1, r_2$, length $\ell$) at center spacing $w$ &
  $S = \dfrac{2\pi \ell}{\operatorname{arcosh}\!\left(\frac{w^2 - r_1^2 - r_2^2}{2r_1 r_2}\right)} \approx \dfrac{2\pi \ell}{\ln\!\left(\frac{w^2}{r_1 r_2}\right)}$ ($w \gg r_1,r_2$) &
  Infinite medium ($k$) between $T_1$ and $T_2$ \\
Eccentric circular cylinders (radii $r_1 < r_2$, length $\ell$, axis offset $e$) &
  $S = \dfrac{2\pi \ell}{\operatorname{arcosh}\!\left(\frac{r_1^2 + r_2^2 - e^2}{2r_1 r_2}\right)}$ &
  Reduces to $\frac{2\pi\ell}{\ln(r_2/r_1)}$ when concentric ($e = 0$) \\
Conduction from thin circular disk (radius $r_0$) on semi-infinite half-space &
  $S = 4 r_0$ \quad (isothermal disk on adiabatic surface); $S = 8 r_0$ (deeply buried) &
  Basis of constriction resistance $R_c = 1/(2 k r_0)$ \\
Uniform volumetric heating $\dot{q}$ $[\si{\watt\per\cubic\meter}]$ in plane slab ($x \in [-L, +L]$) &
  $T(x) - T_s = \dfrac{\dot{q} L^2}{2 k}\!\left(1 - \dfrac{x^2}{L^2}\right) \implies \Delta T_{\max} = \dfrac{\dot{q} L^2}{2 k}$ &
  Symmetrical cooling at $x = \pm L$ (or insulated at $x=0$) \\
Uniform volumetric heating $\dot{q}$ in solid cylinder (outer radius $R$) &
  $T(r) - T_s = \dfrac{\dot{q} R^2}{4 k}\!\left(1 - \dfrac{r^2}{R^2}\right) \implies \Delta T_{\max} = \dfrac{\dot{q} R^2}{4 k}$ &
  Surface flux $q''(R) = \frac{\dot{q}(\pi R^2\ell)}{2\pi R\ell} = \frac{\dot{q} R}{2}$ \\
Uniform volumetric heating $\dot{q}$ in solid sphere (outer radius $R$) &
  $T(r) - T_s = \dfrac{\dot{q} R^2}{6 k}\!\left(1 - \dfrac{r^2}{R^2}\right) \implies \Delta T_{\max} = \dfrac{\dot{q} R^2}{6 k}$ &
  Ratio of $\Delta T_{\max}$ (slab : cyl : sph) $= \frac{1}{2} : \frac{1}{4} : \frac{1}{6} = \mathbf{6 : 3 : 2}$ \\
Maxwell--Eucken dispersed spheres (matrix $k_m$, dispersed pores/particles $k_p$, volume fraction $\phi$) &
  $k_{\mathrm{eff}} = k_m\,\dfrac{2k_m + k_p - 2\phi(k_m - k_p)}{2k_m + k_p + \phi(k_m - k_p)} \xrightarrow{k_p \ll k_m} k_m\,\dfrac{2(1-\phi)}{2+\phi}$ &
  Exact dilute-limit 3-D effective medium between Wiener bounds \\
\bottomrule
\end{tabularx}
\end{minipage}
\par\smallskip

\begin{problem}{P10.7}{Internal Volumetric Joule Heating in Slab vs.\ Cylinder vs.\ Sphere}{10.2}{5.c.iii}{National / MIT}{75 s}{1-D Poisson Equation}{using slab divisor $2k$ instead of cylinder $4k$ or sphere $6k$}
A cylindrical nuclear fuel pin of radius $R = \SI{6.00}{\milli\meter}$ and thermal conductivity $k = \SI{3.00}{\watt\per\meter\per\kelvin}$ generates uniform volumetric thermal power $\dot{q} = \SI{1.50e8}{\watt\per\cubic\meter}$. If the outer cladding interface at $r = R$ is maintained at $T_s = \SI{340.0}{\celsius}$, calculate (a) the outgoing surface heat flux $q''(R)$, (b) the centerline temperature $T(0)$, and (c) what the centerline temperature excess $T(0) - T_s$ would be in a sphere of the same radius $R$, conductivity $k$, and volumetric heating $\dot{q}$.
\end{problem}
\begin{solution}{P10.7}
\begin{enumerate}
  \item[(a)] By steady energy conservation on a cylinder of length $\ell$, all heat generated inside radius $R$ crosses the perimeter $2\pi R \ell$:
  \[
    q''(R) = \frac{\dot{q}\left(\pi R^2 \ell\right)}{2\pi R \ell} = \frac{\dot{q} R}{2} = \frac{(\SI{1.50e8}{\watt\per\cubic\meter})(\SI{0.00600}{\meter})}{2} = \mathbf{\SI{4.50e5}{\watt\per\square\meter}} = \mathbf{\SI{450.0}{\kilo\watt\per\square\meter}}
  \]
  \item[(b)] By \tref{tab:10:shape-factors} (cylindrical Poisson equation $\frac{1}{r}\frac{d}{dr}(r\frac{dT}{dr}) = -\frac{\dot{q}}{k}$), the centerline temperature excess is:
  \[
    \Delta T_{\max,\mathrm{cyl}} = T(0) - T_s = \frac{\dot{q} R^2}{4 k} = \frac{(\SI{1.50e8}{})(0.00600)^2}{4(3.00)} = \frac{5400}{12.0} = \mathbf{\SI{450.0}{\kelvin}} \implies T(0) = 340.0 + 450.0 = \mathbf{\SI{790.0}{\celsius}}
  \]
  \item[(c)] In a sphere of radius $R$, the surface-to-volume ratio is $3/R$ instead of $2/R$, so the divisor is $6k$ instead of $4k$: $\Delta T_{\max,\mathrm{sph}} = \dfrac{\dot{q}R^2}{6k} = \dfrac{2}{3}(\SI{450.0}{\kelvin}) = \mathbf{\SI{300.0}{\kelvin}}$ ($T(0) = \SI{640.0}{\celsius}$).
\end{enumerate}
\end{solution}

\begin{problem}{P10.8}{Buried Hot-Water Pipe Conduction Shape Factor in Soil}{10.4}{5.c.iii}{State / National}{60 s}{Shape Factor $S$}{using $\ln(z/r)$ instead of $\ln(2z/r)$}
A district-heating pipe of outer diameter $D = \SI{10.0}{\centi\meter}$ ($r_1 = \SI{0.0500}{\meter}$) and length $\ell = \SI{25.0}{\meter}$ is buried horizontally with its centerline at depth $z = \SI{0.750}{\meter}$ ($z/r_1 = 15.0$) below a soil surface at $T_{\mathrm{soil}} = \SI{8.0}{\celsius}$. The soil has thermal conductivity $k_{\mathrm{soil}} = \SI{0.950}{\watt\per\meter\per\kelvin}$ (\tref{tab:10:k-hierarchy}) and the pipe outer surface is at $T_p = \SI{78.0}{\celsius}$. Calculate the steady conductive heat loss $\dot{Q}$ to the ground surface.
\end{problem}
\begin{solution}{P10.8}
From \tref{tab:10:shape-factors}, for a buried horizontal cylinder with $z/r_1 = 15.0 > 3$:
\[
  S = \frac{2\pi \ell}{\ln(2z/r_1)} = \frac{2\pi(\SI{25.0}{\meter})}{\ln(30.00)} = \frac{157.080}{3.40120} = \mathbf{\SI{46.184}{\meter}} \implies
  \dot{Q} = S k_{\mathrm{soil}}\left(T_p - T_{\mathrm{soil}}\right) = (46.184)(0.950)(78.0 - 8.0) = \mathbf{\SI{3071}{\watt}} = \mathbf{\SI{3.071}{\kilo\watt}}
\]
\end{solution}

\begin{problem}{P10.9}{Maxwell--Eucken Effective Conductivity of Cellular Cork Insulation}{10.2}{5.c.iii}{National / Invitational}{75 s}{Effective Medium Theory}{using arithmetic parallel average on closed-cell foam}
Closed-cell natural cork consists of a solid suberin cell-wall matrix ($k_m = \SI{0.190}{\watt\per\meter\per\kelvin}$, volume fraction $1 - \phi = 0.220$) enclosing isolated spherical stagnant-air cells ($k_p = \SI{0.0262}{\watt\per\meter\per\kelvin}$, porosity $\phi = 0.780$). Compare (a) the Wiener series lower bound $k_{\mathrm{ser}}$, (b) the Wiener parallel upper bound $k_{\mathrm{par}}$, and (c) the 3-D Maxwell--Eucken dispersed-pore conductivity $k_{\mathrm{ME}}$ against the measured value $k_{\mathrm{cork}} = \SI{0.045}{\watt\per\meter\per\kelvin}$.
\end{problem}
\begin{solution}{P10.9}
\begin{enumerate}
  \item[(a)] \textbf{Wiener Series Lower Bound:} $k_{\mathrm{ser}} = \left(\dfrac{0.220}{0.190} + \dfrac{0.780}{0.0262}\right)^{-1} = (1.1579 + 29.7710)^{-1} = \mathbf{\SI{0.0323}{\watt\per\meter\per\kelvin}}$.
  \item[(b)] \textbf{Wiener Parallel Upper Bound:} $k_{\mathrm{par}} = (0.220)(0.190) + (0.780)(0.0262) = 0.04180 + 0.02044 = \mathbf{\SI{0.0622}{\watt\per\meter\per\kelvin}}$.
  \item[(c)] \textbf{Maxwell--Eucken 3-D Closed-Cell Formula (\tref{tab:10:shape-factors}):}
  \[
    k_{\mathrm{ME}} = (0.190)\,\frac{2(0.190) + 0.0262 - 2(0.780)(0.190 - 0.0262)}{2(0.190) + 0.0262 + (0.780)(0.190 - 0.0262)} = (0.190)\,\frac{0.4062 - 0.25553}{0.4062 + 0.12776} = (0.190)\,\frac{0.15067}{0.53396} = \mathbf{\SI{0.0536}{\watt\per\meter\per\kelvin}}
  \]
  All three satisfy the rigorous theorem $k_{\mathrm{ser}} (\num{0.0323}) < k_{\mathrm{exp}} (\num{0.0450}) < k_{\mathrm{ME}} (\num{0.0536}) < k_{\mathrm{par}} (\num{0.0622})$!
\end{enumerate}
\end{solution}

\begin{problem}{P10.10}{2027 Collector Multi-Wall Conduction Budget \& Fastener Bridge Audit}{10.2}{3.a--h, 5.c.iii}{2027 Device / All Levels}{75 s}{Parallel Conductance Network}{ignoring aluminum corner brackets or floor conduction}
A team's cubic collector enclosure has outer dimensions $\SI{30.0}{\centi\meter}\times\SI{30.0}{\centi\meter}\times\SI{30.0}{\centi\meter}$ with four vertical sidewalls and a bottom floor made of $\SI{35.0}{\milli\meter}$-thick corkboard ($k_{\mathrm{cork}} = \SI{0.0420}{\watt\per\meter\per\kelvin}$, total active conduction area $A_{\mathrm{walls+floor}} = \SI{0.280}{\square\meter}$). The inner cavity air is at $T_{\mathrm{cav}} = \SI{53.0}{\celsius}$ and the room is at $T_{\mathrm{room}} = \SI{21.0}{\celsius}$ ($\Delta T = \SI{32.0}{\kelvin}$). (a) Calculate the steady conductive heat loss $\dot{Q}_{\mathrm{cork}}$ through the corkboard walls and floor. (b) If a student secures the box with $8$ through-wall aluminum machine screws (Alloy 6061-T6, $k_{\mathrm{Al}} = \SI{167}{\watt\per\meter\per\kelvin}$, screw shank diameter $d = \SI{4.00}{\milli\meter}$, length $L = \SI{35.0}{\milli\meter}$), how much extra heat $\dot{Q}_{\mathrm{screws}}$ leaks through the $8$ screws, and what percentage of the total wall conduction do the screws carry?
\end{problem}
\begin{solution}{P10.10}
\begin{enumerate}
  \item[(a)] Through the $\SI{35.0}{\milli\meter}$-thick corkboard enclosure ($A_c = \SI{0.280}{\square\meter}$):
  \[
    G_{\mathrm{cork}} = \frac{k_{\mathrm{cork}} A_c}{L} = \frac{(0.0420)(0.280)}{0.0350} = \mathbf{\SI{0.3360}{\watt\per\kelvin}} \implies
    \dot{Q}_{\mathrm{cork}} = G_{\mathrm{cork}}\Delta T = (0.3360)(32.0) = \mathbf{\SI{10.75}{\watt}}
  \]
  \item[(b)] Each $\SI{4.00}{\milli\meter}$-diameter screw has cross-sectional area $A_1 = \frac{\pi}{4}(0.00400)^2 = \SI{1.2566e-5}{\square\meter}$, so all $8$ screws have total area $A_{\mathrm{screws}} = 8 A_1 = \SI{1.0053e-4}{\square\meter}$ (just $\SI{0.0359}{\percent}$ of the wall area!):
  \[
    G_{\mathrm{screws}} = \frac{k_{\mathrm{Al}} A_{\mathrm{screws}}}{L} = \frac{(167)(\SI{1.0053e-4}{})}{0.0350} = \mathbf{\SI{0.4797}{\watt\per\kelvin}} \implies
    \dot{Q}_{\mathrm{screws}} = (0.4797)(32.0) = \mathbf{\SI{15.35}{\watt}}
  \]
  Eight small screws occupying $\SI{0.036}{\percent}$ of the wall area conduct $\mathbf{\SI{15.35}{\watt}}$ — \textbf{more than the entire corkboard enclosure ($\SI{58.8}{\percent}$ of the combined $\SI{26.10}{\watt}$ leak)}! Always assemble collector walls with non-metallic PVA wood glue, paper tape, or balsa dowels.
\end{enumerate}
\end{solution}

\alsosee{\spref{sec:01.9} (\texttt{QS-09} Conduction Fast Sheet), \spref{sec:11} (Convection \& Dead-Air Gaps), \spref{sec:12} (Radiation Shields), \spref{sec:13} (Transient Lumped Heating \& Biot Number), \spref{sec:14} (Materials Encyclopedia \& \texttt{\#6 PS} Ban), \spref{sec:15} (2027 Collector Cradle Design).}
